\section{半双线性型的判别式}

本小节中，A 表示一个交换环，J 表示 A 的一个自同构，E 表示一个维数为 n 的自由 A-模。

定义 1. --- 给定 E 上关于 J 的半双线性型 \(\Phi\) 以及由 E 的 n 个元素构成的组 \(S = (x_1, \ldots, x_n)\)，称 \(\Phi\) 关于该组的判别式（记作 \(D_{\Phi}(x_1, \ldots, x_n)\) 或 \(D_{\Phi}(S)\)）为 A 的元素 det \((\Phi(x_i, x_j))\)。

若 \((e_{1},\ldots,e_{n})\) 是 E 的一个基，则 \(\Phi\) 关于该基的判别式正是 \(\Phi\) 关于该基的矩阵的行列式。

由 \(\Phi\) 到 \(\wedge\) E 的扩张的定义（\(\S\) 1，第 9 小节）可知

\[
\mathrm{D} _ {\Phi} (x _ {1}, \dots , x _ {n}) = \Phi_ {(n)} (x _ {1} \wedge \dots \wedge x _ {n}, x _ {1} \wedge \dots \wedge x _ {n}),\tag{1}
\]

其中 \(\Phi_{(n)}\) 表示 \(\Phi\) 到 \(\wedge\) E 的扩张。因此，对每个置换 \(\sigma \in \mathfrak{S}_n\)，有

\[
\mathrm{D} _ {\Phi} (x _ {\sigma (\mathbf {1})}, \dots , x _ {\sigma (n)}) = \mathrm{D} _ {\Phi} (x _ {\mathbf {1}}, x _ {\mathbf {2}}, \dots , x _ {n}).
\]

例. --- 设 B 是环 A 上的一个代数，且 B 是维数为 n 的自由 A-模。于是映射 \((x, y) \to \mathrm{Tr}_{\mathrm{B}/\mathrm{A}}(xy)\)（第 VIII 章，§ 12，第 2 小节）是 B 上的一个双线性型。给定由 B 的 n 个元素构成的组 \((x_{1}, \ldots, x_{n})\)，此型关于该组的判别式称为组 \((x_{1}, \ldots, x_{n})\) 在 A 上的判别式，记作 \(\mathrm{D}_{\mathrm{B}/\mathrm{A}}(x_{1}, \ldots, x_{n})\) 。于是有

\[
\mathrm{D} _ {\mathrm{B/A}} (x _ {1}, \dots , x _ {n}) = \det (\operatorname{Tr} _ {\mathrm{B/A}} (x _ {i} x _ {j})).\tag{2}
\]

注. --- 设 \((e_{1},\ldots,e_{n})\) 是 B 在 A 上的一个基，且 \(e_{i}e_{j}=\sum_{k=1}^{n}c_{ijk}e_{k}\;(c_{ijk}\in\mathrm{A})\) 为该基的乘法表（第 II 章，§ 7，第 2 小节）。由于 B 的 A-线性自同态 \(x\to e_{k}x\) 关于 \((e_{r})\) 的矩阵是 \((c_{ksr})\)，故有 \(\mathrm{Tr}_{\mathrm{B}/\mathrm{A}}(e_{k})=\sum_{r=1}^{n}c_{krr}\)，从而 \(\mathrm{Tr}_{\mathrm{B}/\mathrm{A}}(e_{i}e_{j})=\sum_{k,r}c_{ijk}c_{krr}\) 。由此可得

\[
\mathrm{D} _ {\mathbf {B} / \mathbf {A}} (e _ {1}, \dots , e _ {n}) = \det _ {i, j} (\sum_ {k, r} c _ {i j k} c _ {k r r}).\tag{3}
\]

命题 1. — 设 \(\Phi\) 是 E 上关于 J 的半双线性型，\((x_1,\ldots,x_n)\) 是由 E 的 \(n\) 个元素构成的组，\((a_{ij})_{(i,j=1,\ldots,n)}\) 是 A 的 \(n^2\) 个元素构成的族；令 \(y_j = \sum_{i=1}^{n} a_{ji} x_i\) 。则有

\[
\mathrm{D} _ {\Phi} (y _ {1}, \dots , y _ {n}) = \det (a _ {i j}). \det (a _ {i j}) ^ {\mathrm{J}}. \mathrm{D} _ {\Phi} (x _ {1}, \dots , x _ {n}).
\]

事实上，由于 \(\Phi\) 是半双线性型，有 \(\Phi(y_i, y_j) = \sum_{k,m} a_{ik} \Phi(x_k, x_m) a_{jm}^{\mathfrak{J}}\) 。因此，若以 \(A\) 记矩阵 \((a_{ij})\)，则矩阵 \((\Phi(y_i, y_j))\) 等于 \(A \cdot (\Phi(x_i, x_j)) \cdot {}^t A^{\mathfrak{J}}\) 。由 \(\det(^t A) = \det(A)\) 与 \(\det(A^{\mathfrak{J}}) = \det(A)^{\mathfrak{J}}\)，即得我们的断言。

特别地，若 \((e_{i})\) 与 \((e_{i}^{\prime})\) 是 E 的两个基，D 与 D' 是 \(\Phi\) 分别关于这两个基的判别式，a 是从基 \((e_{i})\) 到基 \((e_{i}^{\prime})\) 的过渡矩阵的行列式，则有

\[
\mathrm{D} ^ {\prime} = a a ^ {\mathrm{J}} \mathrm{D}.\tag{4}
\]

由命题 1 可知，若 \((e_i)\) 是 E 的一个基，而 \((x_i)\) 是 E 的 \(n\) 个元素构成的任意一个组，则 \(\mathrm{D}_{\Phi}(e_1,\ldots ,e_n)\) 整除 \(\mathrm{D}_{\Phi}(x_1,\dots,x_n)\) 。特别地，\(\Phi\) 关于 E 的任意两个基的判别式生成 A 的同一个主理想。

设 \((\mathbf{E}_i)_{i\in I}\) 是由有限维自由 A-模构成的有限族，\(\Phi_i\) 是 \(\mathbf{E}_i\) 上关于 J 的半双线性型，\(\mathbf{B}_i\) 是 \(\mathbf{E}_i\) 的一个基。若以 \(\Phi\) 表示诸 \(\Phi_i\) 的直和（\(\S 1\)，第 3 小节），以 B 表示 \(\prod E_i\) 的由诸 \(\mathbf{B}_i\) 之并得到的基，则显然有

\[
\mathrm{D} _ {\Phi} (\mathrm{B}) = \prod_ {i \in \mathrm{I}} \mathrm{D} _ {\Phi_ {i}} (\mathrm{B} _ {i}).\tag{5}
\]

设 \(\Phi\) 是关于 \(J\) 的 \(E\) 上的半双线性型，\(h\) 是 \(A\) 到一个交换环 \(A'\) 中的同态，\(\Phi'\) 是 \(A' \otimes_A E\) 上由 \(\Phi\) 经标量扩张得到的半双线性型（\(\S 1\)，第 4 小节），\((x_1, \ldots, x_n)\) 是 \(E\) 的元素的任意一个组。由于 \(A' \otimes_A E\) 是自由 \(A'\) -模，\(D_{\Phi'}(1 \otimes x_1, \ldots, 1 \otimes x_n)\) 有定义，且显然有

\[
\mathrm{D} _ {\Phi^ {\prime}} (1 \otimes x _ {1}, \dots , 1 \otimes x _ {n}) = h (\mathrm{D} _ {\Phi} (x _ {1}, \dots , x _ {n})).\tag{6}
\]

例. --- 设 B 是 A 上的一个代数，且是维数为 \(n\) 的自由 A-模，\((x_{1},\ldots,x_{n})\) 是 B 在 A 上的一个基，m 是 A 的一个理想。若以 \(h\) 表示 B 到 B/mB 上的典范同态，则 \((h(x_{1}),\ldots,h(x_{n}))\) 是 B/mB 在 A/m 上的一个基（第 I 章，§ 6，第 5 小节，命题 5），且 B/mB 同构于 \((A/m)\otimes_{A}B\) 。于是有

\[
\mathrm{D} _ {(\mathrm{B} / \mathfrak {m B}) / (\mathrm{A} / \mathfrak {m})} (h (x _ {1}), \dots , h (x _ {n})) = h (\mathrm{D} _ {\mathrm{B} / \mathrm{A}} (x _ {1}, \dots , x _ {n})).
\]

命题 2. --- 假设 A 是一个整环。设 \(\Phi\) 是 E 上关于 J 的半双线性型，\((e_1, \ldots, e_n)\) 是 E 的一个基，使得 \(\mathrm{D}_{\Phi}(e_1, e_2, \ldots, e_n) \neq 0\) 。

\begin{enumerate}
\def\labelenumi{\alph{enumi})}
\item
  要使组 \((x_{1},\ldots ,x_{n})\)（由 E 的 \(n\) 个元素组成）是自由的，必须且只需 \(\mathrm{D}_{\Phi}(x_1,\dots,x_n)\) \(\neq 0\) 。
\item
  要使组 \((x_{1},\ldots ,x_{n})\)（由 E 的 \(n\) 个元素组成）是 E 的一个基，必须且只需 \(\mathrm{D}_{\Phi}(x_1,\dots,x_n)\) 与 \(\mathrm{D}_{\Phi}(e_1,\dots,e_n)\) 是 A 中的相伴元素（参见第 VI 章，§ 1，第 5 小节）。
\end{enumerate}

令 \(x_j = \sum_{i=1}^{n} a_{ji} e_i\)（ \(a_{ji} \in A\) ）。先证 \(a\) )。若 \(D_{\Phi}(x_1, \ldots, x_n) = 0\)，则由命题 1 有 \(\det(a_{ji})\) . \(\det(a_{ji})^{\mathfrak{J}} = 0\)，因为 \(D_{\Phi}(e_1, \ldots, e_n) \neq 0\) 且 A 是整环；于是 \(\det(a_{ji}) = 0\)，从而向量 \(x_j\) 线性相关（第 III 章，§ 7，第 1 小节，定理 1，应用于向量空间 \(K \otimes_A E\)，其中 K 表示 A 的分式域）。反之，若这些向量线性相关，则 \(\det(a_{ji}) = 0\)（ \(ibid.\) ），从而 \(D_{\Phi}(x_1, \ldots, x_n) = 0\)（命题 1）。

再证 \(b\) )。若 \(\mathrm{D}_{\Phi}(x_1, \ldots, x_n)\) 与 \(\mathrm{D}_{\Phi}(e_1, \ldots, e_n)\) 在 A 中相伴，则命题 1 表明 \(\det(a_{ij}) \cdot \det(a_{ij})^{\mathrm{J}}\) 在 A 中可逆。于是 \(\det(a_{ij})\) 在 A 中也可逆；从而 A 上的矩阵 \((a_{ij})\) 可逆（第 III 章，§ 6，第 5 小节，定理 2），且 E 的自同态 \(g\)（由 \(g(e_i) = x_i\)（ \(i = 1, \ldots, n\) ）定义）是一个自同构；于是 \((x_1, \ldots, x_n)\) 是 E 的一个基。逆命题由命题 1 立即得出。

命题 3. --- 设 \(\Phi\) 是 E 上关于 J 的半双线性型，S 是 E 的一个基。下列条件等价：

\begin{enumerate}
\def\labelenumi{\alph{enumi})}
\item
  映射 \(s_{\Phi}\)（E 到 \(E^{*}\) 中与 \(\Phi\) 相伴的映射）是双射。
\item
  映射 \(d_{\Phi}\)（E 到 \(E^{*}\) 中与 \(\Phi\) 相伴的映射）是双射。
\item
  元素 \(D_{\Phi}(S)\) 在 A 中可逆。
\end{enumerate}

事实上，条件 \(c\) 表明 \(\Phi\) 关于 S 的矩阵可逆（第 III 章，§ 6，第 5 小节，定理 2）。于是 \(c\) 等价于 \(a\)（ \(§ 1\) ，第 10 小节）；同样，\(c\) 等价于 \(b\)）。

命题 4. --- 假设 A 是一个整环。设 S 是 E 的一个基。要使 E 上的半双线性型 \(\Phi\) 非退化，必须且只需 D \(_{\Phi}\) (S) ≠ 0。

事实上，设 K 是 A 的分式域，\(\Phi'\) 是 \(\Phi\) 到 K-向量空间 \(K\otimes_{\mathbf{A}}E\) 的扩张；我们把 E 与这个向量空间的一个子集视为同一。此时关系 \(D_{\Phi}(S)\neq0\) 由公式 (6) 等价于 \(D_{\Phi'}(S)\neq0\)，而后者表明 \(s_{\Phi'}\) 是双射（命题 3），即 \(\Phi'\) 非退化（\(\S1\)，第 6 小节，命题 6）。又，对每个 \(x\in K\otimes_{\mathbf{A}}E\)，存在 \(a\in A\) 使 \(ax\in E\)。因此，要使 \(\Phi\) 退化，必须且只需 \(\Phi'\) 退化。这就证明了我们的断言。

命题 5. --- 设 A 是一个体，B 是 A 上维数为 n 的交换代数，S 是 B 的一个基。要使 B 可分（第 VIII 章，§ 7，第 5 小节，定义 1），必须且只需 D \(_{B/A}\) (S) ≠ 0。

设 A' 是 A 的代数闭包，B' 是 A' 上的代数 A' ⊗\_A B。若 B 可分，则 B' 半单（第 VIII 章，§ 7，第 5 小节，命题 7 的推论），因而是 n 个同构于 A' 的体的直积（第 VIII 章，§ 6，第 4 小节，命题 9 的推论）。若以 S' 表示 B' 的典范基（与 A'	extsuperscript{n}视为同一），则有 D\_\{B'/A'\}(S') = 1，从而 D\_\{B'/A'\}(S) ≠ 0（命题 1）且 D\_\{B/A\}(S) ≠ 0（公式 (6)）。

反之，设 \(D_{B/A}(S) \neq 0\)。要证明 B 可分，只需证明 \(B'\) 半单，即它不含 \(\neq 0\) 的幂零元素。若 \(x'\) 是 \(B'\) 的一个非零幂零元素，则可取它作为某个基 \(S'\) （\(B'\) 的）的第一个元素，于是 \(\mathrm{Tr}_{\mathbf{B}'/\mathbf{A}'}(x'y') = 0\) 对一切 \(y' \in S'\) 成立，因为幂零自同态的特征值全为零（第 VII 章，§ 5，第 3 小节，命题 8 的推论 3），从而迹为零。由此将得出 \(\mathrm{D}_{\mathbf{B}'/\mathbf{A}'}(\mathrm{S}') = 0\)，于是 \(\mathrm{D}_{\mathbf{B}'/\mathbf{A}'}(\mathrm{S}) = 0\)（命题 1）且 \(\mathrm{D}_{\mathbf{B}/\mathbf{A}}(\mathrm{S}) = 0\)（公式 (6)），与假设矛盾。

注. --- 假设 B 是 A 的一个扩体。设 \(S = (x_{1}, \ldots, x_{n})\) 是 B 的一个基，\((s_{1}, \ldots, s_{n})\) 是 B 到 A 的代数闭包 A' 中的诸 A-同构（每个 \(s_{j}\) 重复 \([B : A]_{i}\) 次）。回忆，对每个 \(z \in B\)，有 \(\operatorname{Tr}_{\mathbf{B}/\mathbf{A}}(z) = \sum_{j=1}^{n} s_{j}(z)\)（第 VIII 章，§ 12，第 2 小节，命题 4）。

于是由行列式的乘法公式可得

\[
(\det (s _ {j} (x _ {i}))) ^ {2} = \mathrm{D} _ {\mathrm{B/A}} (x _ {1}, \dots , x _ {n}).\tag{7}
\]

此公式表明命题 5 推广了第 V 章，§ 7，第 2 小节注中给出的可分性条件。

命题 6. --- 设 \(\Phi\) 是 E 上的一个 A-双线性型，K 是 A 的一个子环，使得 A 是维数为 q 的自由 K-模。若 \((e_i)_{i=1,\ldots,n}\) 是 E 在 A 上的一个基，\((a_j)_{j=1,\ldots,q}\) 是 A 在 K 上的一个基，则 \((a_j e_i)\) 是 E 在 K 上的一个基。映射 \(\Phi'\) 从 E \(\times\) E 到 K 中，由 \(\Phi'(x,y)=\mathrm{Tr}_{\mathrm{A/K}}(\Phi(x,y))\) 定义，它是 E 上的 K-双线性型，且有

\[
\mathrm{D} _ {\Phi^ {\prime}} (a _ {j} e _ {i}) = \mathrm{N} _ {\mathrm{A} / \mathrm{K}} (\mathrm{D} _ {\Phi} (e _ {1}, \dots , e _ {n})) \cdot (\mathrm{D} _ {\mathrm{A} / \mathrm{K}} (a _ {1}, \dots , a _ {q})) ^ {n}.\tag{8}
\]

前两个断言是明显的，只需证明 (7)。按定义，左端是由下式定义的 E 的 K-线性自同态 u 的行列式

\[
u (a _ {j} e _ {i}) = \sum_ {r, s} \mathrm{Tr} _ {\mathbf {A} / \mathbf {K}} (a _ {j} a _ {r} \Phi (e _ {i}, e _ {s})) a _ {r} e _ {s}.
\]

考察 E 的 A-线性自同态 \(\nu\)（由 \(\nu(e_i) = \sum_s \Phi(e_i, e_s)e_s\) 定义）以及 E 的 K-线性自同态 \(w\)（由 \(w(a_j e_i) = (\sum_r \mathrm{Tr}_{\mathrm{A/K}}(a_j a_r)a_r)e_i\) 定义）。我们有

\[
w \left(v \left(a _ {j} e _ {i}\right)\right) = w \left(\sum_ {s} a _ {j} \Phi \left(e _ {i}, e _ {s}\right) e _ {s}\right) = \sum_ {\tau , s} \operatorname{Tr} _ {\mathrm{A} / \mathrm{K}} \left(a _ {j} \Phi \left(e _ {i}, e _ {s}\right) a _ {\tau}\right) a _ {r} e _ {s}
\]

因为 \(w(ae_s) = \sum_r \mathrm{Tr}_{\Lambda/\mathbb{K}}(aa_r)a_r e_s\) 对一切 \(a \in A\) 成立；于是 \(u\) 是复合映射 \(w \circ v\)。这样，以 \(v_{\kappa}\) 表示把 \(v\) 视为 K-线性映射，则有 \(\det(u) = \det(v_{\kappa})\det(w)\)。而 \(\det(v_{\kappa}) = N_{\Lambda/\mathbb{K}}(\det(v))\)（第 VIII 章，§ 12，第 2 小节，命题 7），且显然 \(\det(v) = D_{\Phi}(e_1, \ldots, e_n)\)。另一方面，由于每个 A-模 \(Ae_i\)（ \(i = 1, \ldots, n\) ）在 \(w\) 下稳定，且 \(w\) 在 \(Ae_i\) 上的限制的行列式为 \(\det(\mathrm{Tr}_{\Lambda/\mathbb{K}}(a_j a_r)) = D_{\Lambda/\mathbb{K}}(a_1, \ldots, a_q)\)，故有 \(\det(w) = (D_{\Lambda/\mathbb{K}}(a_1, \ldots, a_q))^n\)。于是公式 (8) 化归为 \(\det(u) = \det(v_{\kappa})\det(w)\)，此式上文已证。

推论（“判别式的传递公式”）. — 设 K 是一个交换环，A 是在 K 上具有有限基 \((a_{j})_{j=1,\ldots,q}\) 的交换代数，E 是在 A 上具有有限基 \((e_{i})_{i=1,\ldots,n}\) 的代数。于是 \((a_{j}e_{i})\) 是 E 在 K 上的一个基，且有

\[
\mathrm{D} _ {\mathrm{E} / \mathrm{K}} (a _ {j} e _ {i}) = \mathrm{N} _ {\mathrm{A} / \mathrm{K}} (\mathrm{D} _ {\mathrm{E} / \mathrm{A}} (e _ {1}, \dots , e _ {n})) \cdot (\mathrm{D} _ {\mathrm{A} / \mathrm{K}} (a _ {1}, \dots , a _ {q})) ^ {n}.
\]

事实上，若令 \(\Phi(x,y)=\mathrm{Tr}_{\mathrm{E/A}}(xy)\)，则由迹的传递公式（第 VIII 章，§ 12，第 2 小节，命题 7 的推论），命题 6 的 K-双线性型 \(\Phi'\) 为 \(\Phi'(x,y)=\mathrm{Tr}_{\mathrm{E/K}}(xy)\)。

习题. — 1) 设 A 是域 K 上有限秩的代数，具有单位元。

\begin{enumerate}
\def\labelenumi{\alph{enumi})}
\item
  证明：若 A 的根不等于零，则 A 上的双线性型 \((x, y) \to \mathrm{Tr}_{\mathbf{A}/\mathbf{K}}(xy)\) 是退化的。
\item
  假设 K 的特征为 0。证明：若 A 是矩阵代数 \(\mathbf{M}_{n}(\mathrm{K})\)，S 是 A 在 K 上的典范基，则有 \(\mathrm{D}_{\mathrm{A}/\mathrm{K}}(\mathrm{S}) \neq 0\)。
\item
  由 a) 与 b) 推出：要使特征 0 的域 K 上有限秩的代数 A 绝对半单，必须且只需双线性型 \((x,y)\to\mathrm{Tr}_{\mathbf{A}/\mathbf{K}}(xy)\) 非退化（或者等价地，对 A 在 K 上的每个基 S 有 \(\mathrm{D}_{\mathbf{A}/\mathbf{K}}(\mathrm{S})\neq0\)）。
\end{enumerate}

设 B 是一个环，A 是 B 的一个包含 B 的单位元的子环；于是 B 是一个 (A, A)-双模；我们以 \(^s\text{B}\)（相应地 \(^d\text{B}\)）表示把集合 B 视为左（相应地右）A-模，以 \(^s\text{B}^*\)（相应地 \(^d\text{B}^*\)）表示 \(^s\text{B}\)（相应地 \(^d\text{B}\)）的右（相应地左）对偶 A-模。对每个 \(x' \in ^s\text{B}^*\) 与每个 \(b \in \text{B}\)，\(x \to \langle xb, x' \rangle\) 是 \(^s\text{B}\) 上的 A-线性型，因而是 \(^s\text{B}^*\) 中记作 \(bx'\) 的元素；映射 \((b, x') \to bx'\) 在 \(^s\text{B}^*\) 上定义了一个左 B-模结构（参见第 III 章，第 2 版，附录 II，第 7 小节）。

\[
(\mathrm{A}, \mathrm{A})
\]

\[
(\mathrm{A}, \mathrm{A}) -
\]

\[
\Phi : (x, y) \rightarrow \varphi (x y)
\]

\(^{s}\) B × \(^{d}\) B 到 A 中的型是非退化的，必须且只需 \(\varphi(0)\) 不包含 B 的除 \(\{0\}\) 之外的任何理想（左理想或右理想）。这时称 \(\varphi\) 为 B 到 A 中的弗罗贝尼乌斯同态。

\begin{enumerate}
\def\labelenumi{\alph{enumi})}
\setcounter{enumi}{1}
\item
  设 \(\varphi\) 是 B 到 A 中的弗罗贝尼乌斯同态；证明映射 \(d_{\Phi}\) （与 \(\Phi\) 右相伴）是左 B-模 \(B_s\) 到左 B-模 \(^sB^*\) 的一个子模上的同构。证明在下列两种情形的每一种中 \(d_{\Phi}\) 都是双射：\(1^\circ\) A 是满足条件（ \(N_s\) ）与（ \(N_d\) ）的左、右阿廷环（\(\S 1\)，习题 11），且 \(^sB\) 与 \(^dB\) 是有限长度的自由 A-模（利用 \(\S 1\) 习题 11 b)）；\(2^\circ\) A 是交换的对合阿廷环（第 VIII 章，\(\S 3\)，习题 11），含于 B 的中心，且 \(^{s}\) B 是有限长度的 A-模（利用第 VIII 章 § 3 习题 11）。
\item
  反之，证明若 B \(_{s}\) 与 \(^{s}\) B \(^{*}\) 同构，则在 b) 所考虑的两种情形之一中，且 \(^{s}\) B 与 \(^{d}\) B 有相同长度时，存在 B 到 A 中的弗罗贝尼乌斯同态（利用 § 1 习题 11 b)）。
\item
  在 a) 的假设与记号下，证明 B 的左理想 I（相应地右理想 r）的右零化子（相应地左零化子）是子模 I（相应地 r）（sB（相应地 dB）的子模）关于型 Φ 的正交补 l⁰（相应地 r⁰）。
\item
  设 \(\varphi\) 是 B 到 A 中的弗罗贝尼乌斯同态。证明若 A 是对合阿廷环（第 VIII 章，§ 3，习题 11），则在再假设 \(b\) ) 的两个条件之一成立时，B 也是（利用 \(b\) ) 与 \(d\) )）。
\end{enumerate}

¶ 3) 设 A 是一个域，B 是 A 上有限秩的代数，具有单位元。

\begin{enumerate}
\def\labelenumi{\alph{enumi})}
\item
  要使 B 是弗罗贝尼乌斯代数，必须且只需存在 B 到 A 中的弗罗贝尼乌斯同态（习题 2）。（利用上文的习题 2 c) 与 e)，以及第 VIII 章 § 13 习题 6 b)。）
\item
  设 \(\varphi\) 是 B 到 A 中的弗罗贝尼乌斯同态；于是 B 上的每个 A-线性型都可以以唯一方式写成 \(x \to \varphi(b'x)\)（相应地 \(x \to \varphi(xb'')\)），其中 \(b'\) 与 \(b''\) 属于 B；要使此型是弗罗贝尼乌斯同态，必须且只需 \(b'\)（相应地 \(b''\)）在 B 中可逆。
\item
  对每个 \(x \in B\)，设 \(x^{\sigma}\) 是满足 \(\varphi(xy) = \varphi(yx^{\sigma})\)（对一切 \(y \in B\)）的唯一元素（参见 b)）。证明 \(x \to x^{\sigma}\) 是 B 的一个 A-自同构。若 \(\sigma\) 是 B 的内自同构，则称弗罗贝尼乌斯代数 B 是对称的；此时存在 B 到 A 中的一个弗罗贝尼乌斯同态，使 \(\sigma\) 为恒等（参见 b)）。这等价于说 (B, B)-双模 B 与 \(^{s}B^{*} = ^{d}B^{*}\)（记作 B*）同构（习题 2 c)）。
\item
  设 E 是有限长度的左 B-模，E' 是它的对偶，E'* 是把 E' 视为 A 上向量空间时的对偶；对 \(x' \in \mathrm{E}', x'' \in \mathrm{E}*\)，\(b \in \mathrm{B}\)，令 \(\langle x', bx'' \rangle = \langle x'b, x'' \rangle\)，便给 E'* 配备了一个左 B-模结构（第 III 章，第 2 版，附录 II，第 7 小节）。对每个 \(x \in \mathrm{E}\)，设 \(f_{\mathrm{E}}(x)\)（或简记 \(f(x)\)）为 E'* 中满足 \(\langle x', f(x) \rangle = \varphi(\langle x, x' \rangle)\)（对一切 \(x' \in \mathrm{E}'\)）的元素；证明 \(f\) 是关于自同构 \(\sigma\) 的半线性双射，从左 B-模 E 映到左 B-模 E'* 上（利用第 VIII 章 § 4 习题 10）。对 \(\mathrm{E} = \mathrm{B}_s\)，有（用习题 2 b) 的记号）\(d_{\Phi}(x^{\sigma}) = f_{\mathrm{B}_s}(x)\)，对一切 \(x \in \mathrm{B}\)。
\end{enumerate}

\begin{enumerate}
\def\labelenumi{\arabic{enumi})}
\setcounter{enumi}{3}
\tightlist
\item
  \begin{enumerate}
  \def\labelenumii{\alph{enumii})}
  \tightlist
  \item
    设 G 是一个有限群。证明：群 G 在任意域 A 上的代数 B 是对称弗罗贝尼乌斯代数（习题 3）。（考察 B 到 A 中的映射 \(\varphi\)，它把每个元素 \(x = \sum_{s \in G} \xi_s.s\) 对应到 \(\varphi(x) = \xi_e\)，其中 \(e\) 表示 G 的单位元。）
  \end{enumerate}
\end{enumerate}

\begin{enumerate}
\def\labelenumi{\alph{enumi})}
\setcounter{enumi}{1}
\tightlist
\item
  设 E 是域 A 上维数为 n 的向量空间，B 是外代数 \(\wedge\) E。证明 B 是弗罗贝尼乌斯代数。（若 \((e_i)_{1 \leqslant i \leqslant n}\) 是 E 的一个基，考察这样的映射：它把 B 的每个元素 \(x\) 对应到 \(e_1 \wedge e_2 \wedge \ldots \wedge e_n\) 在 \(x\) 用与 \((e_i)\) 相应的 B 的基表出时的系数。）要使弗罗贝尼乌斯代数 B 是对称的，必须且只需 \(n\) 是偶数或 A 的特征为 2。
\end{enumerate}

¶ 5) a) 证明：域 K 上有限秩的两个弗罗贝尼乌斯代数（相应地弗罗贝尼乌斯且对称的代数）的张量积是弗罗贝尼乌斯代数（相应地弗罗贝尼乌斯且对称的代数）（参见 § 1，第 9 小节，命题 9）。

\begin{enumerate}
\def\labelenumi{\alph{enumi})}
\setcounter{enumi}{1}
\tightlist
\item
  设 B 是 K 上有限秩的代数，L 是 K 上有限次的扩张。证明：若 L 上的代数 \(\mathrm{B}_{(\mathrm{L})} = \mathrm{B} \otimes_{\mathrm{K}} \mathrm{L}\) 是弗罗贝尼乌斯代数（相应地弗罗贝尼乌斯且对称的代数），则 B 也是。（利用上文的习题 2 c) 与 3 c)，以及第 VIII 章 § 2 习题 2。）
\end{enumerate}

\begin{enumerate}
\def\labelenumi{\arabic{enumi})}
\setcounter{enumi}{5}
\tightlist
\item
  证明：域 A 上有限秩的绝对半单代数 B 都是对称弗罗贝尼乌斯代数。（化归到 B 为单代数的情形；利用上文的习题 5 b)，以及第 VIII 章 § 12 第 3 小节命题 9。）
\end{enumerate}

\section{埃尔米特型与二次型}

本章以下内容中，除另有明确说明外，A 表示一个环，E 表示一个左 A-模。设 A 配备一个对合反自同构 J，记作 \(\alpha \to \overline{\alpha}\)；于是有 \((\overline{\alpha + \beta}) = \overline{\alpha} + \overline{\beta}\)、\((\overline{\alpha\beta}) = \overline{\beta}.\overline{\alpha}\) 及 \(\overline{\overline{\alpha}} = \alpha\)（对 A 中一切 \(\alpha\)、\(\beta\)）。除另有明确说明外，所考虑的半双线性型都是关于该反自同构的右半双线性型（\(\S 1\)，第 2 小节，定义 4）。

\subsection{\texorpdfstring{埃尔米特型与 \(\varepsilon\)-埃尔米特型。}{埃尔米特型与 epsilon-埃尔米特型。}}

定义 1. --- 设 \(\varepsilon\) 是 A 的中心的一个元素。称半双线性型 \(\Phi\)（E 上的）满足 \(\Phi(x, y) = \varepsilon \overline{\Phi(y, x)}\)（对 E 中一切 \(x\) 与 \(y\)）者为 \(\varepsilon\)-埃尔米特型。1-埃尔米特型（相应地 \((-1)\)-埃尔米特型）称为埃尔米特型（相应地反埃尔米特型）。

当 J 为恒等（这蕴含 A 交换）时，（关于 J 的）埃尔米特型（相应地反埃尔米特型）正是对称（相应地反对称）双线性型（第 III 章，§ 5，第 1 小节，定义 2）。回忆，交错双线性型（第 III 章，§ 5，第 2 小节，定义 4）是反对称的；反之，若在 A 中关系 \(2a = 0\) 蕴含 \(a = 0\)，则逆命题成立。

关于 ε-埃尔米特型的正交关系（§ 1，第 3 小节）显然是对称的（参见习题 1）。

若 \(\alpha\) 是 A 的可逆元素，则映射 T: \(\lambda \to \alpha^{-1}\overline{\lambda}\alpha\) 是 A 的一个反自同构，且容易验证型 \(\Phi\alpha\) 关于 T 是半线性的。此外，若 \(\alpha = \overline{\alpha}\)，则 T 是对合的，且若 \(\Phi\) 是 \(\varepsilon\)-埃尔米特型，则 \(\Phi\alpha\) 也是；事实上有

\[
(\lambda^ {\mathrm{T}}) ^ {\mathrm{T}} = \alpha^ {- 1} (\overline {{\alpha^ {- 1} \lambda \alpha}}) \alpha = \alpha^ {- 1} \overline {{\alpha}} \lambda \overline {{\alpha}} ^ {- 1} \alpha = \lambda
\]

\[
\Phi (y, x) \alpha = \varepsilon \overline {{\Phi (x , y)}} \alpha = \varepsilon (\Phi (x, y) \alpha) ^ {\mathrm{T}}.
\]

特别地，当 A 是一个体时，A 的中心的元素 \(\alpha\) 满足 \(\bar{\alpha} = \alpha\) 者组成 A 的一个子域 K，而 E 上（关于 J 的）\(\varepsilon\)-埃尔米特型组成 K 上的一个向量空间。

注. --- 1) 若 \(\Phi\) 是 E 上的 \(\varepsilon\)-埃尔米特型，则 \(\Phi(x, y) = \varepsilon \Phi(x, y) \bar{\varepsilon}\) 对 E 中一切 \(x, y\) 成立。因此，若 \(\Phi\) 取可逆值，则 \(\varepsilon \bar{\varepsilon} = 1\)。

\begin{enumerate}
\def\labelenumi{\arabic{enumi})}
\setcounter{enumi}{1}
\tightlist
\item
  若 A 的中心存在可逆元素 i 使得 \(\bar{i}=i\varepsilon\)，则要使 \(\Phi\) 是 \(\varepsilon\)-埃尔米特型，必须且只需 \(i\Phi\) 是埃尔米特型。
\end{enumerate}

由于映射 \((y, x) \to \overline{\Phi(x, y)}\) 关于 J 是半双线性的，要使 \(\Phi\) 是 \(\varepsilon\)-埃尔米特型，必须且只需有 \(\Phi(y, x) = \varepsilon \overline{\Phi(x, y)}\)（当 \(x\) 与 \(y\) 取遍 E 的一个生成组时）。特别地，若 E 具有有限基 \((e_i)_{1 \leqslant i \leqslant n}\)，则要使 E 上的半双线性型 \(\Phi\) 是 \(\varepsilon\)-埃尔米特型，必须且只需它的矩阵 \(R = (\rho_{ij}) = (\Phi(e_i, e_j))\) 满足关系 \(\rho_{ji} = \varepsilon \overline{\rho_{ij}}\)（对一切 \(i, j\)），即 \(^t R = \varepsilon \overline{R}\)；具有这一性质的矩阵 \(R\) 称为 \(\varepsilon\)-埃尔米特矩阵。当 \(\varepsilon = 1\)（相应地 -1）时，称 \(R\) 关于反自同构 J 是埃尔米特（相应地反埃尔米特）矩阵。当 J 为恒等（从而 A 交换）时，埃尔米特（相应地反埃尔米特）矩阵 \(R\) 满足 \(^t R = R\)（相应地 \(^t R = -R\)）；这时称 \(R\) 为对称（相应地反对称）矩阵。要使 \(\Phi\) 是交错型，必须且只需它的矩阵是反对称的，且此外 R 的对角项全为零；具有这些性质的矩阵称为交错矩阵。

设 \(\Phi\) 是 E 上的一个半双线性型，\(s_{\Phi}\) 与 \(d_{\Phi}\) 是 \(\Phi\) 左相伴与右相伴的 E 到 E* 中的映射（§ 1，第 6 小节）。要使 \(\Phi\) 是 \(\varepsilon\)-埃尔米特型，必须且只需 \(\langle x, s_{\Phi}(y) \rangle = \bar{\varepsilon} \langle x, d_{\Phi}(y) \rangle\) 对 E 的一切元素 \(x, y\) 成立，即 \(s_{\Phi} = \bar{\varepsilon} d_{\Phi}\)，或者等价地，\(\langle x, d_{\Phi}(y) \rangle = \varepsilon \langle x, s_{\Phi}(y) \rangle\)，即 \(d_{\Phi} = \varepsilon s_{\Phi}\)。

设 \(\Phi\) 是一个 \(\varepsilon\)-埃尔米特型，使得映射 \(d_{\Phi}\) （E 到 E* 中的、右相伴于 \(\Phi\) 的）是双射。于是对 E 的每个自同态 \(u\) 有

\[
u ^ {* *} = \varepsilon \bar {\varepsilon} u.\tag{1}
\]

事实上，当 \(x\) 与 \(y\) 是 \(\mathbf{E}\) 中的任意元素时，有

\[
\begin{array}{r l} \Phi (x, u ^ {* *} (y)) & = \Phi (u ^ {*} (x), y) = \varepsilon \overline {{\Phi (y , u ^ {*} (x))}} = \varepsilon \overline {{\Phi (u (y) , x)}} \\ & = \varepsilon \Phi (x, u (y)) \bar {\varepsilon} = \Phi (x, \varepsilon \bar {\varepsilon} u (y)) \end{array}
\]

因此 \(u^{**}(x) = \varepsilon \bar{\varepsilon} u(x)\)，因为 \(\Phi\) 非退化。

若 \(\Phi\) 是 \(\varepsilon\)-埃尔米特型，使得映射 \(s_{\Phi}\) 与 \(d_{\Phi}\) 均为双射，则逆型 \(\hat{\Phi}\)（\(\Phi\) 的逆型）（\(\S 1\)，第 7 小节）是 \(\bar{\varepsilon}\)-埃尔米特型。事实上，为简便起见令 \(s = s_{\Phi}\)、\(d = d_{\Phi}\)，由 \(d = \varepsilon s\) 可得 \(s^{-1} = \bar{\varepsilon} d^{-1}\)（\(s\) 是半线性的）。于是，对 E 中任意的 \(u\)、\(\nu\)，有

\[
\hat {\Phi} (u, \nu) = \Phi (s ^ {- 1} (u), d ^ {- 1} (\nu)) = \bar {\varepsilon} \Phi (d ^ {- 1} (u), d ^ {- 1} (\nu)),
\]

从而

\[
\widehat {\Phi} (\nu , u) = \overline {{\varepsilon \varepsilon}} \overline {{\Phi (d ^ {- 1} (u) , d ^ {- 1} (\nu))}} = \overline {{\varepsilon}} \widehat {\Phi} (u, \nu),
\]

因为 ε 属于 A 的中心。

最后，当环 A 交换时，\(\varepsilon\)-埃尔米特型 \(\Phi\) 到 E 的张量幂与外幂 \(\bigotimes^{p}\mathbf{E}\) 和 \(\bigwedge^{p}\mathbf{E}\) 的典范扩张是 \(\varepsilon^{p}\)-埃尔米特型，这由 § 1 第 9 小节公式 (35) 与 (37) 立即得出。

\subsection{二次扩张上的模。}

设 K 是一个交换环。取 A 为二次扩张 \(\mathrm{A} = \mathrm{K}(i)\)，其中 \(i^{2} = -1\)，取 J 为自同构 \(\lambda + i\mu \to \lambda - i\mu\)（ \(\lambda \in \mathbf{K}\)、\(\mu \in \mathbf{K}\)）（第 II 章，§ 7，第 7 小节）。若 E 是一个 A-模，我们以 \(\mathbf{E}_0\) 表示由 E 经标量环限制得到的 K-模，以 \(j\) 表示自同构 \(x \to ix\)（\(\mathbf{E}_0\) 的）；显然有 \(j^2 = -I\)，其中 I 是 \(\mathbf{E}_0\) 的恒等映射。反之，设 \(\mathbf{E}_0\) 是一个 K-模，\(j\) 是 \(\mathbf{E}_0\) 的满足 \(j^2 = -I\) 的自同构；映射 \(\lambda + i\mu \to \lambda I + \mu j\) 显然是 A 到环 \(\mathfrak{L}(\mathbf{E}_0)\) （\(\mathbf{E}_0\) 的自同态环）中的同态；于是我们就在 \(\mathbf{E}_0\) 上定义了一个 A-模结构，对此有

\[
(\lambda + i \mu) x = \lambda x + \mu j (x) \quad (x \in \mathrm{E} _ {0}, \lambda \in \mathrm{K}, \mu \in \mathrm{K}).\tag{2}
\]

若 E' 是另一个 A-模，\(E_{0}'\) 是 \(E'\) 的底 K-模，\(j'\) 是自同构 \(x' \rightarrow ix'\)（\(E_{0}'\) 的），则 E 到 \(E'\) 中的 A-线性映射 f 正是 \(E_{0}\) 到 \(E_{0}'\) 中的满足 \(f \circ j = j' \circ f\) 的 K-线性映射。特别地，若以 \(E^{*}\) 与 \((\mathrm{E}_{0})^{*}\) 分别表示 E 与 \(E_{0}\) 的对偶，且 \(f_{1}\) 与 \(f_{2}\) 是 E 到 K 中的两个映射，则要使 E 到 A 中的映射 \(x \to f_{1}(x) + if_{2}(x)\) 是 A-线性的，必须且只需 \(f_{1}\) 与 \(f_{2}\) 属于 \((\mathrm{E}_{0})^{*}\)，且 \(f_{1} \circ j + i(f_{2} \circ j) = if_{1} - f_{2}\)，即 \(f_{1} = f_{2} \circ j\) 与 \(f_{1} \circ j = -f_{2}\)。由于 j 是 \(E_{0}\) 的自同构且 \(j^{2} = -I\)，这两个条件等价。消去 \(f_{1}\) 或 \(f_{2}\)，可见公式

\begin{enumerate}
\def\labelenumi{(\arabic{enumi})}
\setcounter{enumi}{2}
\tightlist
\item
\end{enumerate}

\[
f (x) = f _ {1} (x) - i f _ {1} (j (x))\tag{4}
\]

\[
f (x) = f _ {2} (j (x)) + i f _ {2} (x)
\]

\((x \in \mathbf{E}, f \in \mathbf{E}^{*}, f_{1} \in (\mathbf{E}_{0})^{*}, f_{2} \in (\mathbf{E}_{0})^{*})\) 在 \(\mathbf{E}^{*}\) 与 \((\mathbf{E}_{0})^{*}\) 之间建立了两个双射对应。

\subsection{与埃尔米特型相伴的双线性型。}

这里仍假设环 A 是交换环 K 的二次扩张 \(\mathrm{A} = \mathrm{K}(i)\)（其中 \(i^{2} = -1\)），且 J 是 A 的自同构 \(\lambda + i\mu \to \lambda - i\mu\)（ \(\lambda \in K, \mu \in K\)）。设 E 与 \(E'\) 是两个 A-模，\(E_{0}\) 与 \(E_{0}'\) 是 E 与 \(E'\) 的底 K-模，j 与 \(j'\) 是自同构 \(x \to ix\) 与 \(x' \to ix'\)（E 与 \(E'\) 的）（参见第 2 小节）。称 \(E_{0} \times E_{0}'\) 上的 K-双线性型 f 在 j 与 \(j'\) 下不变，如果有

\[
f (j (x), j ^ {\prime} (x ^ {\prime})) = f (x, x ^ {\prime})\tag{5}
\]

对 \(x \in \mathrm{E}_0\) 与 \(x' \in \mathrm{E}_0'\)。把 \(x\) 换成 \(j(x)\)，可见该条件等价于

\[
f (x, j ^ {\prime} (x ^ {\prime})) = - f (j (x), x ^ {\prime})\tag{6}
\]

对一切 \(x \in \mathbf{E}_0\) 与 \(x' \in \mathbf{E}_0'\)。

命题 1. --- 设 \(\Phi_{1}\)（相应地 \(\Phi_{2}\)）是 \(\mathbf{E}_{0} \times \mathbf{E}_{0}^{\prime}\) 上在 \(j\) 与 \(j^{\prime}\) 下不变的 K-双线性型。使 \(\Phi_{1}\)（相应地 \(\Phi_{2}\)）对应于映射 \(\Phi\) （\(\mathbf{E} \times \mathbf{E}^{\prime}\) 到 A 中的、由下式定义的）的那个映射

\begin{enumerate}
\def\labelenumi{(\arabic{enumi})}
\setcounter{enumi}{6}
\tightlist
\item
\end{enumerate}

\[
\Phi (x, x ^ {\prime}) = \Phi_ {1} (x, x ^ {\prime}) + i \Phi_ {1} (x, j ^ {\prime} (x ^ {\prime}))\tag{8}
\]

\[
(\text { resp. } \Phi (x, x ^ {\prime}) = - \Phi_ {2} (x, j ^ {\prime} (x ^ {\prime})) + i \Phi_ {2} (x, x ^ {\prime}))
\]

\((x \in \mathrm{E}, x' \in \mathrm{E})\)，是 \(\mathrm{E}_0 \times \mathrm{E}_0'\) 上在 \(j\) 与 \(j'\) 下不变的 K-双线性型所成的 K-向量空间到 \(\mathrm{E} \times \mathrm{E}'\) 上半双线性型所成的 K-向量空间上的一个同构。再设 \(\mathrm{E} = \mathrm{E}'\)；要使 \(\Phi\) 是埃尔米特型，必须且只需 \(\Phi_1\) 是对称的（相应地 \(\Phi_2\) 是反对称的）（参见习题 4）。

事实上，E × E' 到 A 中的每个映射 \(\Phi\) 都可以以唯一的方式写成 \(\Phi = \Phi_1 + i\Phi_2\) 的形式，其中 \(\Phi_1\) 与 \(\Phi_2\) 是 E × E' 到 K 中的映射。由第 2 小节公式 (3)（相应地 (4)），要使偏映射 \(x \to \Phi(x, x')\) 是 A-线性的，必须且只需 \(\Phi_1\)（相应地 \(\Phi_2\)）关于 \(x\) 是 K-线性的，且有

\begin{enumerate}
\def\labelenumi{(\arabic{enumi})}
\setcounter{enumi}{8}
\tightlist
\item
\end{enumerate}

\[
\Phi (x, x ^ {\prime}) = \Phi_ {1} (x, x ^ {\prime}) - i \Phi_ {1} (j (x), x ^ {\prime})\tag{10}
\]

\[
(\text { resp. } \Phi (x, x ^ {\prime}) = \Phi_ {2} (j (x), x ^ {\prime}) + i \Phi_ {2} (x, x ^ {\prime})).
\]

类似地，要使 \(\overline{\Phi(x,x^{\prime})}\) 关于 \(x^{\prime}\) 是 A-线性的，必须且只需 \(\Phi_{1}\)（相应地 \(\Phi_{2}\)）关于 \(x^{\prime}\) 是 K-线性的，且有

\begin{enumerate}
\def\labelenumi{(\arabic{enumi})}
\setcounter{enumi}{10}
\tightlist
\item
\end{enumerate}

\[
\Phi (x, x ^ {\prime}) = \Phi_ {1} (x, x ^ {\prime}) + i \Phi_ {1} (x, j ^ {\prime} (x ^ {\prime}))\tag{12}
\]

\[
(\text { resp. } \Phi (x, x ^ {\prime}) = - \Phi_ {2} (x, j ^ {\prime} (x ^ {\prime})) + i \Phi_ {2} (x, x ^ {\prime})).
\]

由此立即得出，要使 \(\Phi\) 是半双线性型，必须且只需它可以写成 (9) 与 (11)（相应地 (10) 与 (12)）两种形式之一，且 \(\Phi_{1}\)（相应地 \(\Phi_{2}\)）是在 \(j\) 与 \(j'\) 下不变的 K-双线性型。

由此可知，要使半双线性型 \(\Phi = \Phi_{1} + i\Phi_{2}\) 为零，必须且只需 \(\Phi_{1}\)（相应地 \(\Phi_{2}\)）为零。而若 \(E = E'\)，则 \(\Phi(y, x) = \Phi_{1}(y, x) + i\Phi_{2}(y, x)\) 且 \(\overline{\Phi(x, y)} = \Phi_{1}(x, y) - i\Phi_{2}(x, y)\)；要使这两个表达式相等，换言之要使 \(\Phi\) 是埃尔米特型，必须且只需 \(\Phi_{1}\) 对称（相应地 \(\Phi_{2}\) 反对称）。

注. --- 1) 公式 (7) 与 (8) 表明，若 \(x \in E\)，要使 \(\Phi(x, x') = 0\) 对一切 \(x' \in E'\) 成立，必须且只需 \(\Phi_1(x, x') = 0\)（相应地 \(\Phi_2(x, x') = 0\)）对一切 \(x' \in E'\) 成立。

\begin{enumerate}
\def\labelenumi{\arabic{enumi})}
\setcounter{enumi}{1}
\tightlist
\item
  E 的自同态 \(u\) 关于 \(\Phi\) 的伴随（\(\S 1\)，第 8 小节）与 \(u\)（视为 \(\mathrm{E}_{0}\) 的自同态）关于 \(\Phi_{1}\)（相应地 \(\Phi_{2}\)）的伴随相同。
\end{enumerate}

\subsection{二次型。}

定义 2. --- 假设环 A 交换。称 E 到 A 中的映射 Q 为 E 上的二次型，如果

\begin{enumerate}
\def\labelenumi{\arabic{enumi})}
\item
  对 α ∈ A 与 x ∈ E 有 Q(αx) = α²Q(x)；
\item
  映射 \(\Phi : (x, y) \to \mathrm{Q}(x + y) - \mathrm{Q}(x) - \mathrm{Q}(y)\) （\(\mathbf{E} \times \mathbf{E}\) 到 A 中的）是一个双线性型。
\end{enumerate}

双线性型 \(\Phi\)（它必然是对称的）称为与 Q 相伴的双线性型。若 \(\Phi\) 非退化，则称 Q 非退化。

由于 Q(2x) = 4Q(x)，由 2) 立即得出

\[
\Phi (x, x) = 2 \mathrm{Q} (x).\tag{13}
\]

特别地，若 A 是特征为 2 的环，则型 \(\Phi\) 是交错型。

我们称 E 的两个元素（相应地两个子集）关于 Q 正交，如果它们关于相伴的双线性型 Φ 正交。

设 \((x_{i})_{i\in I}\) 是 E 的元素的一个族，\((a_{i})_{i\in I}\) 是 A 的元素的一个族，其中除有限多个外均为零。对指标 \(i\)（使 \(a_{i}\neq0\) 者）的个数作归纳，容易证明有

\[
\mathrm{Q} (\sum_ {i} a _ {i} x _ {i}) = \sum_ {i} a _ {i} ^ {2} \mathrm{Q} (x _ {i}) + \sum_ {\{i, j \}} a _ {i} a _ {j} \Phi (x _ {i}, x _ {j}),\tag{14}
\]

最后一个和式遍取 I 的二元子集。

对每个双线性型 \(f\) （在 \(\mathbf{E} \times \mathbf{E}\) 上），定义二次型 \(\mathbf{Q}\) 为 \(\mathrm{Q}(x) = f(x, x)\)；于是双线性型 \(\Phi\) （与 \(\mathbf{Q}\) 相伴的）由 \(\Phi(x, y) = f(x, y) + f(y, x)\)（对 \(x, y\)，在 \(\mathbf{E}\) 中）定义。此外，若假设标量 2 有逆 \(\frac{1}{2}\) （在 \(\mathbf{A}\) 中），则存在唯一的对称双线性型 \(f\) 使 \(\mathrm{Q}(x) = f(x, x)\)，即 \(f = \frac{1}{2}\Phi\)；\(f\) 关于组 \(S = (x_1, \ldots, x_n)\) 的判别式也称为 \(\mathbf{Q}\) 关于 \(S\) 的判别式。因此在这种情形，\(\mathbf{E}\) 上的二次型与对称双线性型之间存在一一对应（参见习题 6）。

在自由模的情形，我们还有如下结果：

命题 2. --- 假设 A 交换且 E 具有基 \((e_i)_{i \in I}\)。于是，对 E 上的每个二次型 Q，存在一个双线性型 \(f\) （E \(\times\) E 上的），使对一切 \(x \in E\) 有 Q(x) = f(x, x)。对每个标量族 \((b_{ij})_{(i,j) \in I \times I}\)，它满足 \(b_{ij} = b_{ji}\)（对 \((i, j) \in I \times I\)），则存在唯一的二次型 Q 使

\[
\mathrm{Q} (e _ {i}) = b _ {i i}, \quad \Phi (e _ {i}, e _ {j}) = b _ {i j} \text {   for   } i \neq j,\tag{15}
\]

其中 Φ 表示与 Q 相伴的双线性型；于是 Q 由下列公式给出

\[
\mathrm{Q} (\sum_ {i} a _ {i} e _ {i}) = \sum_ {\{i, j \}} b _ {i j} a _ {i} a _ {j},\tag{16}
\]

最后一个和式遍取 I 的具有一个或两个元素的子集 \(\{i, j\}\)。

事实上，由于公式 (16) 只是公式 (14) 的一个改写，满足 (15) 的二次型 Q 的唯一性得证。为证其存在性，首先注意存在 A 的元素的一个族（ \(b_{ij}'\) ），使 \(b_{ii}' = b_{ii}\) 且 \(b_{ij}' + b_{ji}' = b_{ij}\)（对 \(i \neq j\)）；例如，给 I 配备一个全序结构（集合论，第 III 章，§ 2，第 3 小节，定理 1），并令 \(b_{ij}' = b_{ij}\)（对 \(i<j\)），令 \(b_{ij}^{\prime}=0\)（对 \(i>j\)），便得到这样一个族。由于诸 \(e_{i}\) 组成 E 的基，存在一个双线性型 \(f\) （\(\mathbf{E}\times\mathbf{E}\) 上的）使 \(f(e_{i},e_{j})=b_{ij}^{\prime}\)；令 \(\mathrm{Q}^{\prime}(x)=f(x,x)\)，并以 \(\Phi^{\prime}\) 表示与二次型 \(\mathrm{Q}^{\prime}\) 相伴的双线性型，则得 \(\mathrm{Q}^{\prime}(e_{i})=b_{ii}\) 且 \(\Phi^{\prime}(e_{i},e_{j})=f(e_{i},e_{j})+f(e_{j},e_{i})=b_{ij}\)。这就证明了我们的第二个断言。至于第一个断言，它立即得出，因为据唯一性，若二次型 Q 满足 (15)，则有 \(\mathrm{Q}(x)=\mathrm{Q}^{\prime}(x)=f(x,x)\)。

配备了由二次型 Q 所定义的结构的模 E 称为二次模。二次模 (E, Q) 到二次模 (E', Q') 中的同态是指一个 E 到 E' 中的线性映射 u，使 Q = Q' ∘ u；若 Φ 与 Φ' 是与 Q 和 Q' 相伴的双线性型，则对 x ∈ E, y ∈ E 有 Φ(x, y) = Φ'(u(x), u(y))；换言之，Φ' 是 Φ 在 u 下的原像（§ 1，第 1 小节）。两个 A-模 E 与 E' 上的两个二次型 Q 与 Q' 称为等价的，如果相应的二次模同构。

设 \((\mathbf{E}_{i}, \mathbf{Q}_{i})_{i\in\mathbb{I}}\) 是二次模的一个族，E 是诸模 \(E_{i}\) 的直和。称二次模 \((\mathbf{E}_{i}, \mathbf{Q}_{i})\) 的外直和为这样一个二次模：它是在 E 上配备由 \(\mathrm{Q}(\sum_{i}x_{i}) = \sum_{i}\mathrm{Q}_{i}(x_{i})\)（对 \(x_{i} \in E_{i}\)）定义的二次型 Q 而得到的。也称二次型 Q 为诸二次型 \(Q_{i}\) 的外直和。

若诸型 \(Q_{i}\) 非退化，则 Q 也非退化。

设 Q 是 A-模 E 上的一个二次型；若 F 是 E 的一个子模，且 Q 在模 F 的每个类上取常值，则由 Q 取商得到的 E/F 到 A 中的映射 \(\overline{Q}\) 显然是一个二次型，且 E 到 E/F 上的典范映射是二次模结构的同态。要使 Q 在模 F 的每个类上取常值，必须且只需 \(Q(x + y) = Q(x)\)（对 \(x \in E\) 与 \(y \in F\)），即（以 \(\Phi\) 表示与 Q 相伴的双线性型）\(Q(y) + \Phi(x, y) = 0\)（对 \(y \in F\) 与 \(x \in E\)）。取 x = 0，可见 \(Q(y) = 0\)（对 \(y \in F\)），从而 \(\Phi(x,y)=0\)（对 \(x\in E\) 与 \(y\in F\)）。换言之，若把二次模 (E, Q) 的核定义为 E 中满足 \(Q(x)=0\) 且 \(\Phi(x,z)=0\)（对一切 \(z\in E\)）的元素 x 的集合 N，则要使 Q 在模 F 的每个类上取常值，必须且只需 F 含于 (E, Q) 的核 N。不难验证 N 是 E 的一个子模。因此，要使 Q 在模 F 的每个类上取常值，必须且只需 F 由 N 的元素生成。

立即可见二次模 E/N 的核是 \{0\}。

命题 3. --- 设 h 是 A 到交换环 A' 中的一个同态。对 A-模 E 上的每个二次型 Q，存在唯一的 A'-模 A'⊗A E（第 III 章，第 2 版，附录 II，第 10 小节）上的二次型 Q'，使得

\[
\mathrm{Q} ^ {\prime} (1 \otimes x) = h (\mathrm{Q} (x))\tag{17}
\]

对一切 \(x \in \mathbf{E}\)。此外，与 Q' 相伴的双线性型 \(\Phi'\) 是由与 Q 相伴的双线性型 \(\Phi\) 经标量环扩张得到的。

我们先证明，若存在满足 (17) 的二次型 Q'，则它是唯一的，且与 Q' 相伴的双线性型 \(\Phi'\) 是由与 Q 相伴的型 \(\Phi\) 经标量环扩张得到的。事实上，后一断言由以下事实得出

\[
\begin{array}{r l} \Phi^ {\prime} (1 \otimes x, 1 \otimes y) & = Q ^ {\prime} (1 \otimes x + 1 \otimes y) - Q ^ {\prime} (1 \otimes x) - Q ^ {\prime} (1 \otimes y) \\ & = h (\Phi (x, y)) \end{array}\tag{18}
\]

对 \(x \in \mathrm{E}, y \in \mathrm{E}\)。公式 (14) 于是表明有

\[
\mathrm{Q} ^ {\prime} \left(\sum_ {i} a _ {i} ^ {\prime} \otimes x _ {i}\right) = \sum_ {i} a _ {i} ^ {\prime 2} h (\mathrm{Q} (x _ {i})) + \sum_ {\{i, j \}} a _ {i} ^ {\prime} a _ {j} ^ {\prime} h (\Phi (x _ {i}, x _ {j}))\tag{19}
\]

对 \(a_{i}^{\prime}\in\mathbf{A}^{\prime}\) 与 \(x_{i}\in\mathbf{E}\)，这就证明了 \(\mathbf{Q}^{\prime}\) 的唯一性。

为证 Q' 的存在性，先假设模 E 具有基 \((e_{i})_{i\in I}\)。于是存在 A 的元素 \(b_{ij}\)，使 \(b_{ij}=b_{ji}\) 且 \(\mathrm{Q}(\sum_{i}a_{i}e_{i})=\sum_{\{i,j\}}b_{ij}a_{i}a_{j}\) 对 \(a_{i}\in A\) 成立（命题 2）。由于诸元素 \(1\otimes_{A}e_{i}\) 组成

A'-模 A' ⊗\_A E 的基，故在该模上取下式便定义了该模上的一个二次型 Q'

\[
\mathrm{Q} ^ {\prime} (\sum_ {i} a _ {i} ^ {\prime} \otimes e _ {i}) = \sum_ {\{i, j \}} a _ {i} ^ {\prime} a _ {j} ^ {\prime} h (b _ {i j})\tag{20}
\]

其中 \(a_i' \in A'\)；于是，对 \(x = \sum_{i} a_i e_i \in E\)

\[
Q ^ {\prime} (1 \otimes x) = Q ^ {\prime} (\sum_ {i} h (a _ {i}) \otimes e _ {i}) = \sum_ {\{i, j \}} h (a _ {i}) h (a _ {j}) h (b _ {i j}) = h (Q (x)),\tag{21}
\]

这就证明了这种情形下 Q' 的存在性。

现在转到一般情形。设 \((x_{i})_{i\in I}\) 是 E 的一个生成组，\(\mathbf{A}^{(1)}\) 是 I 的元素的形式线性组合模（第 II 章，§ 1，第 8 小节），\((e_{i})_{i\in I}\) 是 \(\mathbf{A}^{(1)}\) 的典范基。线性映射 \(u\) （\(\mathbf{A}^{(1)}\) 到 E 中的、由 \(u(e_{i}) = x_{i}\) 定义的）是满射，因为诸元素 \(x_{i}\) 生成 E。由此（第 III 章，第 2 版，附录 II，第 5 小节，命题 4）可知映射 \(1\otimes u\) （\(\mathbf{A}'\otimes\mathbf{A}^{(1)}\) 到 \(\mathbf{A}'\otimes\mathbf{E}\) 中的）是满射，且它的核 P' 由形如 \(1\otimes p\) 且 \(u(p)=0\) 的元素生成。于是设 \(\mathrm{Q}_{1}^{\prime}\) 是 \(\mathbf{A}'\otimes_{\mathbf{A}}\mathbf{A}^{(1)}\) 上把二次型 \(\mathrm{Q}_{1}=\mathrm{Q}\circ u\) （\(\mathbf{A}^{(1)}\) 上的）加以扩张得到的二次型。若 \(p\) 是 \(\mathbf{A}^{(1)}\) 中满足 \(u(p)=0\) 的元素，则有 \(\mathrm{Q}_{1}^{\prime}(1\otimes p)=h(\mathrm{Q}_{1}(p))=0\)，且（以 \(\Phi_{1}^{\prime}\) 表示与 \(\mathrm{Q}_{1}^{\prime}\) 相伴的双线性型）\(\Phi_{1}^{\prime}(1\otimes p,1\otimes x)=h(\Phi(u(p),u(x)))=0\)（对一切 \(x\in\mathbf{A}^{(1)}\)）。因此，若 \(u(p)=0\)（ \(p\in\mathbf{A}^{(1)}\) ），则 \(1\otimes p\) 属于二次模 \(\mathbf{A}'\otimes_{\mathbf{A}}\mathbf{A}^{(1)}\) 的核，从而如上所述，存在 \(\mathbf{A}'\otimes_{\mathbf{A}}\mathbf{E}\) 上的一个二次型 Q' 使 \(\mathrm{Q}_{1}^{\prime}=\mathrm{Q}^{\prime}\circ(1\otimes u)\)。由于 \(u\) 是满射，可见 Q' 满足条件 (17)。证毕。

其存在性与唯一性由命题 3 保证的二次型 Q′ 称为 Q 到 A′ ⊗\_A E 的（关于 h 的）扩张。也称 Q' 是由 Q 经标量环扩张得到的。

习题. — 1) 设 A 是一个体，E 是 A 上的一个左向量空间，\(\Phi\) 是 E 上的一个半双线性型（关于 A 的某个反自同构 J）。假设 \(\Phi\) 的秩（有限或无限）\(\geqslant 2\)，且关系 \(\Phi(x, y) = 0\) 与 \(\Phi(y, x) = 0\) 等价。

\begin{enumerate}
\def\labelenumi{\alph{enumi})}
\item
  证明：存在 A 中的 \(\lambda\neq0\) 使 \(\Phi(y,x)=\lambda(\Phi(x,y))^{\mathfrak{J}}\)（利用 § 1 习题 8）。
\item
  证明：存在 \(\alpha\in A\)，使半双线性型 \(\Phi\alpha\)（关于反自同构 \(\xi\to\alpha^{-1}\xi^{J}\alpha\)）是埃尔米特型或交错型。
\end{enumerate}

（首先注意有 \(\xi^{J^2} = \lambda^{-1}\xi\lambda^{-J}\) 且 \(\lambda\lambda^J = \lambda^J\lambda = 1\)。然后区分两种情形：是 \(\xi + \xi^J\lambda^{-1} = 0\) 对一切 \(\xi \in A\) 成立，还是不然；在第二种情形，证明每个 \(\neq 0\) 的形如 \(\alpha = \xi + \xi^J\lambda^{-1}\) 的元素都满足要求。）

\begin{enumerate}
\def\labelenumi{\arabic{enumi})}
\setcounter{enumi}{1}
\tightlist
\item
  设 \(\Phi\) 是体 A 上有限维向量空间 E 上的 \(\varepsilon\)-埃尔米特半双线性型。
\end{enumerate}

\begin{enumerate}
\def\labelenumi{\alph{enumi})}
\item
  证明：对 E 的每个向量子空间 M，有 \((\mathbf{M}^{\mathbf{0}})^{\mathbf{0}} = \mathbf{M} + \mathbf{E}^{\mathbf{0}}\) 且 \(\dim M^{\mathbf{0}} + \dim M = \dim E + \dim (M \cap E^{\mathbf{0}})\)。
\item
  若 \(M_{1}\)、\(M_{2}\) 是 E 的两个向量子空间，证明有 dim \((\mathbf{M}_{1} \cap \mathbf{M}_{2}^{0}) + \dim (\mathbf{M}_{2} + \mathbf{M}_{1}^{0}) = \dim \mathbf{E} + \dim (\mathbf{M}_{1} \cap \mathbf{E}^{0})\)（考察 E 到 \(E/E^{0}\) 的典范映射）。
\end{enumerate}

\begin{enumerate}
\def\labelenumi{\arabic{enumi})}
\setcounter{enumi}{2}
\tightlist
\item
  设 K 是一个交换环，\(f\) 是 K{[}X{]} 中的一个首一多项式，次数为 \(n \geqslant 1\)；设 A 是商代数 K{[}X{]}/(f)，它以 1, \(\xi, \xi^2, \ldots, \xi^{n-1}\) 为 K 上的基（第 IV 章，§ 1，第 5 小节，命题 4）。证明：一个 A-模 E 的数据等价于一个 K-模 \(E_0\) 与 K-自同态 \(j\) （\(E_0\) 的、满足 \(f(j) = 0\) 的）的数据。对每个
\end{enumerate}

\(u \in \mathrm{E}^*\)，令 \(u(x) = \sum_{k=0} u_k(x)\xi^k\)；证明：若 \(\alpha_0 = f(0)\) 在 K 中可逆，则映射 \(u \to u_0\) 是 \(\mathbf{E}^*\) 到 \((\mathbf{E}_0)^*\) 上的一个 K-同构，并写出其逆同构。

¶ 4) a) 设 A 是一个环（交换或不交换），σ 是 A 的一个自同构，使得存在可逆元素 γ ∈ A 满足 γ\^{}σ = γ，且对一切 ξ ∈ A 有 ξ\^{}σ² = γξγ\^{}-1。设 B 是具有两个元素的基 (e₁, e₂) 的左 A-模；证明：在 B 中取下式作为乘法，便在 B 上定义了一个环结构

\[
(\xi e _ {1} + \eta e _ {2}) (\xi^ {\prime} e _ {1} + \eta^ {\prime} e _ {2}) = (\xi \xi^ {\prime} + \eta \eta^ {\prime \sigma} \gamma) e _ {1} + (\eta \xi^ {\prime \sigma} + \xi \eta^ {\prime}) e _ {2}.
\]

对这个环结构，\(e_1\) 是单位元（与 A 的单位元 1 视为同一）；若令 \(e_2 = \rho\)，则 \(\rho^2 = \gamma\) 且 \(\rho\xi = \xi^\sigma\rho\)（对一切 \(\xi \in A\)）。此外，B 是一个右 A-模，1 与 \(\rho\) 组成它的基。若 A 是一个体，则 B 是体的必要充分条件是 \(\gamma\) 不是形如 \(\lambda^\sigma\lambda\)（其中 \(\lambda \in A\)）的元素。（参见第 VIII 章，§ 12，习题 8。）

\begin{enumerate}
\def\labelenumi{\alph{enumi})}
\setcounter{enumi}{1}
\tightlist
\item
  设 J 是 A 的一个对合反自同构。假设存在可逆元素 \(\delta\in\mathsf{A}\) 满足下列条件：
\end{enumerate}

\[
\delta^ {J} = \delta , \quad \delta^ {\sigma} \delta = \gamma \gamma^ {J}, \quad (\xi^ {J}) ^ {\sigma} = \delta (\xi^ {\sigma}) ^ {J} \delta^ {- 1} \quad \text { for   all } \quad \xi \in A.\tag{1}
\]

证明：这时可把 J 扩张为 B 的一个对合反自同构（仍记作 J），办法是令

\[
(\xi + \eta \rho) ^ {J} = \xi^ {J} + \gamma^ {- 1} \delta^ {\sigma} (\eta^ {J}) ^ {\sigma} \rho .\tag{2}
\]

此外，若 A 与 B 都是体，且 \(\sigma\) 不是内自同构，证明条件 (1) 是使 J 可扩张为 B 的对合反自同构的必要条件（令 \(\rho^{J} = \alpha + \beta\rho\)，其中 \(\alpha, \beta\) 属于 A，证明必有 \(\alpha = 0\)，其法是写出条件 \((\rho\xi)^{J} = \xi^{J}\rho^{J}\)（对一切 \(\xi \in A\)））。

\begin{enumerate}
\def\labelenumi{\alph{enumi})}
\setcounter{enumi}{2}
\tightlist
\item
  假设条件 (1) 成立，且对合反自同构 J 已由 (2) 扩张到 B 上。设 F 是一个单式 B-模，E 是把 F 的标量环限制到 A 而得到的单式 A-模；映射 \(j: x \to \rho x\) 于是是 E 到其自身上的一个半线性双射，它关于 A 的自同构 \(\sigma\)，且满足 \(j^{2}(x) = \gamma x\)。设 \(\Phi\) 是 F 上的一个埃尔米特型（关于 J）；对 \(x \in E\)、\(y \in E\)，令 \(\Phi(x, y) = \Phi_{1}(x, y) + \Phi_{2}(x, y)\rho\)，其中 \(\Phi_{1}(x, y) \in A\)、\(\Phi_{2}(x, y) \in A\)。证明 \(\Phi_{1}\) 是 E 上的一个埃尔米特型（关于 J），且满足
\end{enumerate}

\[
\Phi_ {1} (j (x), j (y)) = (\Phi_ {1} (x, y)) ^ {\sigma} \delta ;
\]

我们有 \(\Phi_{2}(x,y)=\Phi_{1}(x,j(y))\delta^{-1}\)，\(\Phi_{2}\) 是 E 上关于反自同构（一般非对合）\(\xi\to(\xi^{\mathrm{J}})^{\sigma}\) 的半双线性型，且满足

\[
\Phi_ {2} (j (x), j (y)) = (\Phi_ {2} (x, y)) ^ {\sigma} \delta^ {\sigma}
\]

and

\[
\Phi_ {2} (y, x) = \gamma^ {\mathrm{J}} \delta^ {- 1} ((\Phi_ {2} (x, y)) ^ {\mathrm{J}}) ^ {\sigma}.
\]

逆命题。要使 \(\Phi\) 非退化，必须且只需 \(\Phi_{1}\) 或 \(\Phi_{2}\) 非退化。B 是交换环 K 上对应于 K 的元素对 \((\alpha, \beta)\)（其中 \(\beta\) 可逆）的四元数代数、A 是 B 的子代数 K + Ku、J 与 \(\sigma\) 是映射 \(\xi \to \bar{\xi}\)（第 II 章，§ 7，第 8 小节）的特殊情形。

\begin{enumerate}
\def\labelenumi{\alph{enumi})}
\setcounter{enumi}{3}
\tightlist
\item
  设 u 是 A-模 E 的一个自同构。要使 u 是 B-模 F 的保持型 \(\Phi\) 不变的自同构，必须且只需 u 满足下列三个条件中的任意两个：
\end{enumerate}

\(1^{\circ} u\) 保持 \(\Phi_{1}\) 不变；

\(2^{o} u\) 保持 \(\Phi_{2}\) 不变；

\(3^{o} u\) 与 \(j\) 交换。

\begin{enumerate}
\def\labelenumi{\arabic{enumi})}
\setcounter{enumi}{4}
\tightlist
\item
  设 A 是一个交换环，E 是一个 A-模，\((x_{i})_{i\in I}\) 是 E 的一个生成组；设 R 是模 \(\mathrm{A}^{(\mathrm{I})}\) 中由元素 \((y_{i})_{i\in\mathrm{I}}\)（满足 \(\sum_{i}y_{i}x_{i}=0\)）组成的子模，\((a_{\lambda})_{\lambda\in\mathrm{L}}\) 是 R 的一个生成组（其中 \(a_{\lambda}=(a_{\lambda i})_{i\in\mathrm{I}}\)）。设 \((b_{ij})\) 是 A 的元素的一个族 \((i\in\mathrm{I},j\in\mathrm{I})\)。要存在一个二次型 Q 使 \(\mathrm{Q}(x_{i})=b_{ii}\) 且 \(\Phi(x_{i},x_{j})=b_{ij}\)（对 \(i\neq j\)）（ \(\Phi\) 表示与 Q 相伴的双线性型），必须且只需对 I 中一切 i, j 有 \(b_{ij}=b_{ji}\)，且对一切 \(\lambda\in L\) 与 \(i\in I\) 有
\end{enumerate}

\[
\sum_ {\{i, j \}} b _ {i j} a _ {\lambda i} a _ {\lambda j} = \sum_ {j \neq i} b _ {i j} a _ {\lambda j} + 2 b _ {i i} a _ {\lambda i} = 0;\tag{3}
\]

于是有 \(\mathrm{Q}(\sum_{i}a_{i}x_{i})=\sum_{\{i,j\}}b_{ij}a_{i}a_{j}\)。由此推出第 4 小节命题 3 的一个新证明。（注意 \(x_{i}^{\prime}=1\otimes x_{i}\) 组成 \(A^{\prime}\otimes_{A}E\) 的一个生成组，且 \(A^{\prime}\)-模 \(A^{\prime}\otimes_{A}E\) 同构于 \(A^{\prime(I)}/R^{\prime}\)，其中 \(A^{\prime(I)}\) 与 \(A^{\prime}\otimes_{A}A^{(I)}\) 视为同一，而 \(R^{\prime}\) 由 R 在 \(A^{(I)}\) 到 \(A^{\prime(I)}\) 中的典范映射下的像生成。）

\begin{enumerate}
\def\labelenumi{\arabic{enumi})}
\setcounter{enumi}{5}
\tightlist
\item
  设 A 是特征 2 的交换环，E 是自由 A-模，\(\mathcal{A}\)（相应地 \(\mathcal{S},\mathcal{Q}\)）是 E 上交错（相应地对称双线性、二次）型所成的 A-模。我们有 \(\mathcal{A}\subset\mathcal{S}\)；再定义线性映射 \(\omega\) （\(\mathcal{S}\) 到 \(\mathcal{Q}\) 中的）与线性映射 \(\theta\) （\(\mathcal{Q}\) 到 \(\mathcal{A}\) 中的）如下：对每个双线性型 \(\Phi\in\mathcal{S}\)，\(\omega(\Phi)\) 是二次型 \(x\to\Phi(x,x)\)；对每个二次型 \(Q\in\mathcal{Q}\)，\(\theta(Q)\) 是与 Q 相伴的双线性型，它是交错的。证明 \(\omega(0)=\mathcal{A}\)、\(\theta(\mathcal{Q})=\mathcal{A}\) 且 \(\bar{\theta}(0)=\omega(\mathcal{S})\)。
\end{enumerate}

¶ 7) 设 A 是一个交换环，E、F 是两个 A-模。称 E 到 F 中的映射 Q 为二次映射，如果它满足下列条件：\(1^{\circ} \mathrm{Q}(\alpha x) = \alpha^{2} \mathrm{Q}(x)\)（对 \(\alpha \in \mathrm{A}\)、\(x \in \mathrm{E}\)）；\(2^{\circ}\) 映射 \((x, y) \to \mathrm{Q}(x + y) - \mathrm{Q}(x) - \mathrm{Q}(y)\) （\(\mathrm{E} \times \mathrm{E}\) 到 F 中的）是双线性的。若 \(f\) 是 A-模 \(\mathrm{E}_{1}\) 到 E 中的线性映射，则 \(\mathrm{Q} \circ f\) 是 \(\mathrm{E}_{1}\) 到 F 中的二次映射。

\begin{enumerate}
\def\labelenumi{\alph{enumi})}
\tightlist
\item
  设 E 是一个 A-模，\(\mathrm{A}^{(\mathrm{E})}\) 是 E 的元素以 A 为系数的形式线性组合模（第 II 章，§ 1，第 8 小节），且对每个 \(x \in \mathrm{E}\)，设 \(\varepsilon_x\) 是 \(\mathrm{A}^{(\mathrm{E})}\) 的典范基中相应的元素。设 \(\Gamma^2(\mathrm{E})\) 是 \(\mathrm{A}^{(\mathrm{E})} \times (\mathrm{E} \otimes_{\mathbf{A}} \mathrm{E})\) 关于由形如 ( \(\varepsilon_{x+y} - \varepsilon_x - \varepsilon_y, -x \otimes y\) ) 与 ( \(\varepsilon_{\lambda x} - \lambda^2\varepsilon_x, 0\) )（对 \(x \in \mathrm{E}\)、\(y \in \mathrm{E}\)、\(\lambda \in \mathrm{A}\)）的元素生成的子模 R 的商。对每个 \(x \in \mathrm{E}\)，令 \(\gamma(x) = \varphi(\varepsilon_x, 0)\)，其中 \(\varphi\) 表示 \(\mathrm{A}^{(\mathrm{E})} \times (\mathrm{E} \otimes \mathrm{E})\) 到 \(\Gamma^2(\mathrm{E})\) 上的典范映射；称 \(\gamma\) 为 E 到 \(\Gamma^2(\mathrm{E})\) 中的典范映射。证明 \(\gamma\) 是 E 到 \(\Gamma^2(\mathrm{E})\) 的二次映射，且对每个 E 到 A-模 F 中的二次映射 Q，存在唯一的线性映射 \(q\) （\(\Gamma^2(\mathrm{E})\) 到 F 中的）使 \(Q = q \circ \gamma\)（换言之，\((\Gamma^2(\mathrm{E}), \gamma)\) 是一个泛映射问题的解；参见集合论，第 IV 章，§ 3）。
\end{enumerate}

对每对 A-模 E, E′ 以及 E 到 E′ 中的每个线性映射 f，证明：若 γ′ 表示 E′ 到 Γ²(E′) 中的典范映射，则存在唯一的 Γ²(E) 到 Γ²(E') 中的线性映射 \(\bar{f}\) 使 γ'∘f = \(\bar{f} \circ \gamma\)。

\begin{enumerate}
\def\labelenumi{\alph{enumi})}
\setcounter{enumi}{1}
\item
  假设 E 是两个子模 M, N 的直和。定义 \(\Gamma^{2}(\mathrm{E})\) 到模 \(\Gamma^{2}(\mathrm{M})\)、\(\Gamma^{2}(\mathrm{N})\) 与 \(\mathbf{M} \otimes \mathbf{N}\) 的直和上的一个典范同构（证明此直和是与 \(\Gamma^{2}(\mathrm{E})\) 相同的泛映射问题的解）。
\item
  设 F 是 E 的一个子模，j 是 F 到 E 中的典范单射。定义 \(\Gamma^{2}(\mathrm{E}/\mathrm{F})\) 到
\end{enumerate}

\[
\Gamma^ {2} (\mathrm{E}) / (\bar {j} (\Gamma^ {2} (\mathrm{F})) + \psi (\mathrm{E} \times \mathrm{F})),
\]

其中 \(\psi(x, y) = \varphi(0, x \otimes j(y))\)（对 \(x \in \mathrm{E}, y \in \mathrm{F}\)）。（同一方法。）

\begin{enumerate}
\def\labelenumi{\alph{enumi})}
\setcounter{enumi}{3}
\item
  设 A' 是一个交换环，h 是 A 到 A' 中的一个同态。定义 \(\Gamma^{2}(\mathbf{A}'\otimes_{\mathbf{A}}\mathbf{E})\) 到 \(\mathbf{A}'\otimes_{\mathbf{A}}\symbf{\Gamma}^{2}(\mathbf{E})\) 上的一个典范同构（同一方法）。
\item
  存在唯一的线性映射 \(s\) （\(\Gamma^{2}(\mathrm{E})\) 到 \(\mathrm{E}\otimes\mathrm{E}\) 中的），使 \(s(\gamma(x))=x\otimes x\) 对一切 \(x\in\mathrm{E}\) 成立；证明若 \(\mathrm{E}\) 是自由模，则 s 是到 E 上 2 阶对称张量所成的子模上的一个同构。
\item
  假设 A = Z 且 E 是 n 阶有限循环群。证明：若 n 为奇数，\(\Gamma^{2}(E)\) 是 n 阶循环群；若 n 为偶数，则是 2n 阶循环群。（首先注意，若 a 是 E 的一个生成元，则 \(\gamma(a)\) 是 \(\Gamma^{2}(E)\) 的生成元，且对每个整数 h 有 \(\gamma(-ha) = \gamma(ha)\)；由此推出，若 n 为奇数，则 \(n\gamma(a) = 0\)（取 \(h = (n - 1)/2\)）；同样证明若 n 为偶数则 \(2n\gamma(a) = 0\)。最后证明：若 n 为奇数（相应地偶数），存在 E 到 n（相应地 2n）阶循环群中的把 a 映到该群生成元上的二次映射 Q。）
\end{enumerate}

\begin{enumerate}
\def\labelenumi{\arabic{enumi})}
\setcounter{enumi}{7}
\tightlist
\item
  设 A 是一个域，E、F 是 A 上的两个向量空间。设 g 是 E 到 F 中的一个映射，使得存在 E × E 到 F 中的三个映射 a, b, c，满足恒等式
\end{enumerate}

\[
g (\lambda x + \mu y) = \lambda^ {2} a (x, y) + \lambda \mu b (x, y) + \mu^ {2} c (x, y)\tag{4}
\]

对 E 中一切 x, y 与 A 中 λ, μ 成立。

\begin{enumerate}
\def\labelenumi{\alph{enumi})}
\tightlist
\item
  证明有 \(a(x, y) = g(x)\)、\(c(x, y) = g(y)\)、\(b(y, x) = b(x, y)\) 及 \(b(\lambda x, y) = \lambda b(x, y)\)；此外，若令
\end{enumerate}

\[
d (x, y, z) = b (x + y, z) - b (x, z) - b (y, z)
\]

证明有 \(d(x,y,z)=d(y,z,x)=d(z,x,y)\)，并由此推出 \(d(\lambda x,\mu y,\nu z)=\lambda\mu\nu d(x,y,z)\)。

\begin{enumerate}
\def\labelenumi{\alph{enumi})}
\setcounter{enumi}{1}
\tightlist
\item
  由 a) 推出：若 \(\mathbf{A} \neq \mathbf{F}_2\)，则必有 \(d(x, y, z) = 0\)，从而 \(g\) 是二次映射（习题 7）。反之，若 \(\mathbf{A} = \mathbf{F}_2\) 且 dim \(\mathrm{E} \geqslant 3\)，证明存在 E 到 A 中的映射 \(g\)，满足 (4) 且使 \(d(x, y, z)\) 不恒为零。
\end{enumerate}

\begin{enumerate}
\def\labelenumi{\arabic{enumi})}
\setcounter{enumi}{8}
\tightlist
\item
  设 A 是一个交换环，E 是具有 n 个元素的基的 A-模，Φ 是 E 上的对称双线性型。设 e 是 \(\bigwedge^{n}\) E 中组成该模的基的一个元素，\(\Delta\) 是 \(\Phi\) 关于 \(e\) 的判别式，\(\varphi_{p}\) 是 \(\bigwedge^{p}\) E 到 \(\bigwedge^{n-p}\) E* 上的典范同构（关于 \(e\)）（第 III 章，§8，第 5 小节）。对 \(x \in \bigwedge^{p}\) E，令 \(d_{(p)}(x) = \bar{\varphi}_{n-p}^{-1}(d_{\Phi(p)}(x))\)；证明映射 \(d_{(p)}\) （\(\bigwedge^{p}\) E 到 \(\bigwedge^{n-p}\) E 中的）具有下列性质：
\end{enumerate}

\begin{enumerate}
\def\labelenumi{\alph{enumi})}
\item
  对每个 \(x \in \wedge \mathrm{E}\)，有 \(d_{(n - p)}(d_{(p)}(x)) = (-1)^{p(n - p)}\Delta x\)。
\item
  对每对元素 \(x, y\) （属于 \(\wedge\) E），有
\end{enumerate}

\[
x \wedge d _ {(p)} (y) = \Phi_ {(p)} (x, y) e \quad \text { and } \quad \Phi_ {(n - p)} (d _ {(p)} (x), d _ {(p)} (y)) = \Delta \Phi_ {(p)} (x, y).
\]

\begin{enumerate}
\def\labelenumi{\alph{enumi})}
\setcounter{enumi}{2}
\item
  进一步假设 A 是一个体且 \(\Phi\) 非退化。那么，若 \(x\) 是对应于 E 的子空间 F 的一个可分解 \(p\)-向量 \(\neq 0\)（第 III 章，§ 7，第 3 小节），则 \(d_{(p)}(x)\) 是一个可分解 \((n-p)\)-向量 \(\neq 0\)，对应于与 F 正交的子空间 F \(^{0}\)。
\item
  将前面的结果推广到如下情形：（A 为交换环，）\(\Phi\) 是 \(\varepsilon\)-埃尔米特半双线性型（关于 A 的一个对合自同构 \(J \neq 1\)）。
\end{enumerate}

\begin{enumerate}
\def\labelenumi{\arabic{enumi})}
\setcounter{enumi}{9}
\tightlist
\item
  \begin{enumerate}
  \def\labelenumii{\alph{enumii})}
  \tightlist
  \item
    设 A 是一个交换环，E 是具有 3 个元素的基的 A-模，\(\Phi\) 是 E 上的对称双线性型。沿用习题 9 的记号，对 E 的任意两个元素 \(x\)、\(y\)，令 \(x \overline{\wedge} y = d_{(2)}(x \wedge y)\)，并称此元素为向量积
  \end{enumerate}
\end{enumerate}

\[
\stackrel {{3}} {\wedge} \mathrm{E})
\]

即 \(x\) 与 \(y\) 的向量积（关于 \(\Phi\) 以及 \(e\) （\(\wedge E\) 的基））。证明 \((x, y) \to x \overline{\wedge} y\) 是 \(E \times E\) 到 \(E\) 中的交错双线性映射，且 \(x \overline{\wedge} y\) 与 \(x\) 以及 \(y\) 都正交。

\begin{enumerate}
\def\labelenumi{\alph{enumi})}
\setcounter{enumi}{1}
\tightlist
\item
  设 \(\alpha, \beta\) 是 A 的两个可逆元素，B 是与偶 \((\alpha, \beta)\) 对应的 A 上的四元数代数（第 II 章，§ 7，第 8 小节），1, u, v, w 是 B 在 A 上的典范基；设 E 是以 u, v, w 为基的 B 的子模。证明：若 x, y 是属于 E 的两个四元数，则有
\end{enumerate}

\[
x y = \Phi (x, y) + x \overline {{{{\wedge}}}} y
\]

其中 \(\Phi\) 是 E 上的对称双线性型，与 \(\Phi\) 相伴的线性映射是双射，而 \(x\overline{\wedge}y\) 是 x 与 y 关于型 \(\Phi\) 以及基 \(\alpha^{-1}\beta^{-1}u\wedge v\wedge w\) （\(\bigwedge^{3}\mathbf{E}\) 的）的向量积。

\begin{enumerate}
\def\labelenumi{\arabic{enumi})}
\setcounter{enumi}{10}
\tightlist
\item
  设 \(\Phi\) 是有限维向量空间 E 上的非退化 \(\varepsilon\)-埃尔米特半双线性型。若两个向量子空间 M、N° 中的一个包含另一个，则称 E 的向量子空间 M （关于 \(\Phi\)）与向量子空间 N 弱正交。
\end{enumerate}

\begin{enumerate}
\def\labelenumi{\alph{enumi})}
\item
  证明关系“M 与 N 弱正交”是对称的。
\item
  若 M 与 N 弱正交，则 \(M^{0}\) 与 \(N^{0}\) 弱正交。
\item
  若 M 与 N 弱正交且 M ∩ N = \{0\}，则 M 与 N 正交。
\end{enumerate}

\section{全迷向子空间。维特定理}

在本节中，除非另有明确说明，均假设 A 是一个体。用 \(\Phi\) 表示 E 上的一个 \(\varepsilon\)-埃尔米特型（关于 A 的对合反自同构 \(\lambda \to \overline{\lambda}\)），或者与 E 上的二次型 Q 相伴的对称双线性型（在后一情形假设 A 是交换的）。

\subsection{迷向子空间。}

定义 1. --- 给定环 A 上的模 E，若 \(\Phi(x, x) = 0\)，则称 E 的元素 x 是迷向的。称 E 的子模 F 是

\begin{enumerate}
\def\labelenumi{\arabic{enumi})}
\item
  迷向的，如果存在 F 中与 F 正交的元素 \(x \neq 0\)；
\item
  全迷向的，如果 \(\Phi\) 在 F 上的限制为零。
\end{enumerate}

当模 E 装备了二次型 Q 时，若 E 的元素关于与 Q 相伴的双线性型是迷向的，则称该元素是迷向的（相应地，若 E 的子模关于与 Q 相伴的双线性型是迷向的或全迷向的，则称该子模是迷向的或全迷向的）。

迷向向量就是与自身正交的向量。说子模 F 是迷向的，就意味着 \(F \cap F^{0} \neq \{0\}\)，或者等价地，\(\Phi\) 在 F 上的限制是退化的；因此 E 的非迷向子模 G 是使 \(\Phi\) 在 G 上的限制非退化的子模。要使 E 的子模 F 是全迷向的，必须且只需 \(F \subset F^{0}\)。若 F 是 E 的全迷向子模，则包含在 F 中的每个子模 \(F'\) 也是全迷向的。两两正交的全迷向子模族之和是全迷向子模。E 的全迷向子模的集合按包含关系排序显然是归纳的；由此推出，每个全迷向子模都包含在某个极大全迷向子模之中。

命题 1. --- 假设 A 是一个体。设 F 是 E 的有限维非迷向子空间；则 E 是 F 与 F⁰ 的直和。

事实上，由于按假设 \(\Phi'\) （\(\Phi\) 在 F 上的限制）非退化，故映射 \(d_{\Phi'}\) （F 到其对偶 F* 的、右相伴于 \(\Phi'\) 的）是单射，又因 F 与 F* 是相同有限维的两个空间，故它是双射。于是，对每个 \(y \in E\)，存在 F 中唯一的元素 \(y_0\)，使得 \(\Phi(x, y) = \Phi(x, y_0)\) 对所有 \(x \in F\) 成立，亦即 \(y - y_0 \in F^0\)；这就证明了 E 是 F 与 F \(^0\) 的直和。

推论. --- 若 F 是 E 的有限维子空间，且 Φ 非退化，则下列条件等价：

\begin{enumerate}
\def\labelenumi{\alph{enumi})}
\item
  F 非迷向。
\item
  \(\mathbf{F}^0\) 非迷向。
\item
  E 是 F 与 F \(^{0}\) 的直和。
\end{enumerate}

命题 1 确实表明 \(a)\) 蕴含 \(c)\)，而 \(c)\) 蕴含 \(a)\) 与 \(b)\)。最后，若 \(\mathbf{F}^0\) 非迷向，则有 \(\mathbf{F} \cap \mathbf{F}^0 \subset \mathbf{F}^0 \cap \mathbf{F}^{00} = \{0\}\)；于是 \(\mathbf{F}\) 非迷向，这表明 \(b)\) 蕴含 \(a)\)。

定义 2. --- 设 Q 是 E 上的二次型。若 Q(x) = 0，则称 E 的元素 x 是奇异的（关于 Q 的）。称 E 的子模 F 是：

\begin{enumerate}
\def\labelenumi{\arabic{enumi})}
\item
  奇异的，如果存在 F 中奇异且与 F 正交的元素 \(x \neq 0\)；
\item
  全奇异的，如果 Q 在 F 上的限制为零。
\end{enumerate}

二次模 (E, Q) 的核（§ 3，第 4 小节）由 E° 的奇异元素组成；要使子模 F 是奇异的，必须且只需其核 ≠ \{0\}。由于 Φ(x, y) = Q(x + y) − Q(x) − Q(y)，每个 ≠ \{0\} 的全奇异子模都是奇异的。由于 Φ(x, x) = 2Q(x)，每个奇异向量都是迷向的，且每个奇异（相应地，全奇异）子模都是迷向的（相应地，全迷向的）；当 A 是特征 ≠ 2 的体时，逆命题成立。包含在全奇异子模中的每个子模本身都是全奇异的。两两正交的全奇异子模族之和是全奇异子模。E 的全奇异子模的集合按包含关系排序是归纳的；因此 E 的每个全奇异子模都包含在某个极大全奇异子模之中。

\subsection{维特分解。}

除了本小节开头以来已经生效的约定之外，我们再补充如下约定：

条件 (T). --- 对所有 \(x \in E\)，存在 \(\alpha \in A\) 使得 \(\Phi(x, x) = \alpha + \varepsilon\bar{\alpha}\)。

当 \(\Phi\) 是交错型，或当 \(\varepsilon = 1\) 且 A 是特征 \(\neq 2\) 的体时，取 \(\alpha = \frac{1}{2}\Phi(x, x)\)，该条件总是满足的（参见习题 1 与 14）。

引理 1. --- 设 \(\Phi\) 是 E 上满足 (T) 的 \(\varepsilon\)-埃尔米特型（相应地，与二次型 Q 相伴的双线性型），F 是 E 的不约化为 0 的全迷向（相应地，全奇异）子空间。对每个不与 F 正交的 \(x \in E\) 和所有 \(\alpha \in A\)，存在 \(y \in F\) 使得

\[
\Phi (x + y, x + y) = \alpha + \varepsilon \bar {\alpha} \quad (\text { resp. } Q (x + y) = \alpha).
\]

事实上，设 \(\Phi(x, x) = \beta + \varepsilon\bar{\beta}\)（相应地，\(Q(x) = \beta\)）。那么对 \(y \in F\)，有 \(\Phi(x + y, x + y) = (\beta + \Phi(x, y)) + \varepsilon(\overline{\beta + \Phi(x, y)})\)（由于 \(\Phi(y, y) = 0\)）（相应地，有 \(Q(x + y) = \beta + \Phi(x, y)\)，由于 \(Q(y) = 0\)）。由于 \(x\) 不与 \(F\) 正交，仿射线性函数 \(y \to \Phi(x, y) + \beta\) （\(F\) 上的）不是常数；因此它取值 \(\alpha\)（对某个元素 \(y\)，该元素属于 \(F\)），于是即回答了问题。

E 的维特分解是指将 E 分解为三个子空间 F、F'、G 的直和，使得 F 与 F' 全迷向（相应地，全奇异）且 G 非迷向并与 F + F' 正交；若 E 是有限维的，则 Φ 关于适应于 E 的维特分解的 E 的基的矩阵具有形式

\[
\left( \begin{array}{c c c} 0 & U & 0 \\ \varepsilon \overline {{U}} & 0 & 0 \\ 0 & 0 & V \end{array} \right)\tag{1}
\]

我们称 \(\Phi\) 为中性型，如果它非退化，且 E 是有限维的并且是两个全迷向（相应地，全奇异）子空间的直和。两个中性型的直和是中性型。

命题 2. --- 设 Φ 是满足 (T) 的非退化 ε-埃尔米特型（相应地，与非退化二次型 Q 相伴的双线性型），F 是有限维 r 的全迷向（相应地，全奇异）子空间。

\begin{enumerate}
\def\labelenumi{\alph{enumi})}
\item
  若 \(F'\) 是维数为 r 且满足 \(F' \cap F^{0} = \{0\}\) 的全迷向子空间，则 \(F + F'\) 非迷向，且对 F 的任意基 \((f_{i})\)，存在基 \((f_{i}')\) （\(F'\) 的），使得 \(\Phi(f_{i}, f_{j}') = \delta_{ij}\)（克罗内克指标）对 \(i, j = 1, \ldots, r\) 成立。
\item
  若 G 是维数 \(\leqslant r\) 且满足 \(G \cap F^{0} = \{0\}\) 的全迷向（相应地，全奇异）子空间，则存在全迷向（相应地，全奇异）子空间 \(F' \supset G\)，维数为 r 且满足 \(F' \cap F^{0} = \{0\}\)。
\end{enumerate}

设 \(\Psi\) 是 \(\Phi\) 在 \(F \times F'\) 上的限制；对 \(x' \in F'\)，关系 \(\langle \Phi(x, x') = 0\) 对所有 \(x \in F\rangle\) 蕴含 \(x = 0\)，因为 \(F' \cap F^0 = \{0\}\)。于是断言 \(a)\) 由 §1 第 6 小节命题 6 的推论得出，只是 \(F + F'\) 非迷向这一点除外。而子空间 \(H = (F + F') \cap (F + F')^0\) 等于 \((F + F') \cap F^0 \cap F'^0\)。由于 \(F \subset F^0\)，有 \((F + F') \cap F^0 = F + (F' \cap F^0) = F\)，从而 \(H = F \cap F'^0\)；故 \(H = \{0\}\)（既然我们已看到 \(\Psi\) 非退化）。这就证明了 \(F + F'\) 非迷向。

为证明 \(b\) )，我们对 \(s = \dim G\) 作递降归纳。因此只需证明：若 \(s < r\)，则存在全迷向（相应地，全奇异）子空间 \(G'\)，它包含 \(G\)、维数为 \(s + 1\) 且满足 \(G' \cap F^0 = \{0\}\)。由于 \(\dim G < \dim F\)，\(\Phi\) 在 \(F \times G\) 上的限制是退化的，又因 \(G \cap F^0\) 为零，故 \(F \cap G^0\) 非零。若此时有 \(G + F^0 \supset G^0\)，则取正交子空间并注意 \(F = F^{00}\) 与 \(G = G^{00}\)（ \(\S 1\) ，第 6 小节，命题 4 的推论 1），便可推出 \(G^0 \cap F \subset G\)，从而

\[
\mathrm{G} ^ {0} \cap \mathrm{F} \subset \mathrm{G} \cap \mathrm{F} \subset \mathrm{G} \cap \mathrm{F} ^ {0} = \{0 \},
\]

这是不可能的。于是存在元素 \(x\) （\(G^0\) 的）使得 \(x \notin G + F^0\)；由于 \(F \subset F^0\)，可以给 \(x\) 加上 \(G^0 \cap F\) 的一个向量而不改变这些性质；由于 \(G^0 \cap F\) 全迷向（相应地，全奇异）且 \(\neq \{0\}\)，引理 1 表明可以选择 \(x\) 是迷向的（相应地，奇异的）。

于是子空间 \(G' = G + Ax\) 的维数为 \(s + 1\)，且是全迷向的（相应地，全奇异的）；此外我们有 \(G' \cap F^{0} = \{0\}\)，因为若 \(y = z + ax\)（\(z \in G, a \in A\)）属于 \(F^{0}\)，则有 a = 0，否则将有 \(x \in F^{0} + G\)，与 x 的选取相矛盾，从而 \(y = z \in G \cap F^{0} = \{0\}\)，即 y = 0。因此子空间 \(G'\) 回答了问题。

推论 1. --- 若 F 是维数为 r 的全迷向（相应地，全奇异）子空间，则存在维数为 r 的全迷向（相应地，全奇异）子空间 F'，使得 F ∩ F' = \{0\} 且 F + F' 非迷向。

在命题 2 b) 中取 G = \{0\} 即可。

推论 2. --- A 上相同维数的空间上的两个中性 ε-埃尔米特型是等价的。

注. --- 在推论 1 的条件下，E 是 F + F' 与 F + F' 的正交子空间的直和。这样就得到 E 的一个维特分解。由命题 2 a)，存在 F 与 F' 的基，使得在 Φ 的矩阵 (1) 中，块 U 是单位矩阵 1。

命题 3. --- 设 \(\Phi\) 是满足 (T) 的非退化 \(\varepsilon\)-埃尔米特型（相应地，与非退化二次型 Q 相伴的双线性型）。设 \(F_{1}\) 与 \(F_{2}\) 是 E 的两个极大全迷向（相应地，全奇异）子空间，其中一个是有限维的。令 \(F = F_{1} \cap F_{2}\)。设 \(S_{i}\)（\(i = 1, 2\)）是 F 在 \(F_{i}\) 中的补；令 \(S = S_{1} + S_{2}\)。则存在 E 的两个子空间 G 和 H，使得

\begin{enumerate}
\def\labelenumi{\alph{enumi})}
\item
  子空间 G + F、S 与 H 非迷向且两两正交；
\item
  E 是 F、S、G 与 H 的直和；
\item
  H 中没有非零的迷向（相应地，奇异）向量；
\item
  G 是全迷向的（相应地，全奇异的）。
\end{enumerate}

此外，\(F_{1}\) 与 \(F_{2}\) 都是有限维的，并且有 \(\dim \mathrm{F}_1 = \dim \mathrm{F}_2, \dim \mathrm{G} = \dim \mathrm{F}, \dim \mathrm{S}_1 = \dim \mathrm{S}_2, \operatorname{codim} \mathrm{H} = 2 \dim \mathrm{F}_1.\)

首先注意，若 N 是极大全迷向（相应地，全奇异）子空间，则每个与 N 正交的迷向（相应地，奇异）向量 x 都属于 N，因为否则 \(N + Ax\) 将与 N 的极大性矛盾。于是，若对 i = 1, 2，\(x_{i}\) 是 \(F_{i}^{0}\) 的迷向（相应地，奇异）向量，则有 \(x_{i} \in F_{i}\)。另一方面，若 y 是 \(S_{1}\) 中与 \(S_{2}\) 正交的元素，则由于 \(F_{1}\) 全迷向，y 与 \(F_{1}\) 正交，从而与 F 正交，进而与 \(F_{2} = S_{2} + F\) 正交。由于 y 是迷向（相应地，奇异）的且与 \(F_{2}\) 正交，故有

\[
y \in S _ {1} \cap F _ {2} = S _ {1} \cap F _ {1} \cap F _ {2} = S _ {1} \cap F = \{0 \}.
\]

于是有 \(S_1 \cap S_2^0 = \{0\}\)，同样 \(S_2 \cap S_1^0 = \{0\}\)。由于两个子空间 \(F_1, F_2\) 之一（例如 \(F_1\)）是有限维的，故 \(S_1\) 有限维，从而 \(S_1^0\) 有有限余维数（§ 1，第 6 小节，命题 4 的推论 1），于是 \(S_2\) 也是有限维的，因为 \(S_2 \cap S_1^0 = \{0\}\)；此外这表明 dim \(S_2 \leqslant \text{codim } S_1^0 = \text{dim } S_1\)；同样 dim \(S_1 \leqslant \text{dim } S_2\)，从而 dim \(S_1 = \text{dim } S_2\)。于是命题 2 a) 表明 \(S = S_1 + S_2\) 非迷向。

在此基础上，S 的正交子空间 N 非迷向（第 1 小节，命题 1 的推论）且包含 F；因此命题 2 的推论 1 表明，存在 N 的全迷向（相应地，全奇异）子空间 G，使得 dim G = dim F、G ∩ F = \{0\} 且 G + F 非迷向。于是 G 满足 d)。再取 H 为 G + F 在 N 中的正交子空间，便满足 a) 与 b)。至于 c)，注意由于 H 与 F₁ = S₁ + F 正交，根据证明开头所作的说明以及 H ∩ F₁ = \{0\} 这一事实，H 中没有非零的迷向（相应地，奇异）向量。最后，关于维数的某些断言已在证明过程中得出；其余的由它们平凡地推出。

推论 1. --- 在命题 3 的假设下，两个有限维的极大全迷向（相应地，全奇异）子空间有相同的维数。对每个有限维的极大全迷向（相应地，全奇异）子空间 F，存在另一个这样的子空间 \(\mathbf{F}'\)，使得 \(\mathbf{F} \cap \mathbf{F}' = \{0\}\)，且在这些条件下 \(\mathbf{F} + \mathbf{F}'\) 非迷向。

若 \(F \cap F' = \{0\}\)，则按命题 3 的记号将有 \(G = \{0\}\)，且 \(F + F'\) 非迷向。其余断言由命题 3 与命题 2 的推论 1 平凡地得出。

推论 2. --- 设 Q 是代数闭体 A 上有限维 n 的向量空间 E 上的非退化二次型；则存在 E 的基 \((e_i)_{1 \leqslant i \leqslant n}\)，使得

\begin{enumerate}
\def\labelenumi{(\arabic{enumi})}
\setcounter{enumi}{1}
\tightlist
\item
\end{enumerate}

\[
\mathrm{Q} (\sum_ {i = 1} ^ {n} x _ {i} e _ {i}) = \sum_ {i = 1} ^ {\nu} x _ {i} x _ {i + \nu} \quad \text {   if   } n = 2 \nu ,\tag{3}
\]

\[
\mathrm{Q} \left(\sum_ {i = 1} ^ {n} x _ {i} e _ {i}\right) = \sum_ {i = 1} ^ {\nu} x _ {i} x _ {i + \nu} + x _ {2 \nu + 1} ^ {2} \quad \text {   if   } n = 2 \nu + 1.
\]

事实上，设 \(F_{1}\) 与 \(F_{2}\) 是满足 \(F_{1} \cap F_{2} = \{0\}\)（推论 1）的两个极大全奇异子空间，q 是它们的维数。于是按命题 3 的记号有 \(G = \{0\}\)。取基 \((e_{i})_{1 \leqslant i \leqslant q}\) （\(F_{1}\) 的）与基 \((e_{i})_{q+1 \leqslant i \leqslant 2q}\) （\(F_{2}\) 的），后者满足 \(\Phi(e_{i}, e_{j+q}) = \delta_{ij}\)（对 \(i, j = 1, \ldots, q\)）（命题 2 a)），我们看出只需证明 dim \(H \leqslant 1\)。而若 \(x \in H\)、\(y \in H\) 且 \(x \neq 0\)，则方程 \(Q(y - ax) = Q(y) - a\Phi(x, y) + a^{2}Q(x) = 0\) 至少有一个解 \(a_{0}\)（由于 \(Q(x) \neq 0\)），又由于 H 的每个奇异向量都是零，故有 \(y = a_{0}x\)。

定义 3. --- 假设 E 是有限维的，且 \(\Phi\) 是满足 (T) 的非退化 \(\varepsilon\)-埃尔米特型（相应地，与非退化二次型 Q 相伴的双线性型）。\(\Phi\)（相应地，Q）的指数是 E 的极大全迷向（相应地，全奇异）子空间的公共维数。

若 n 是 E 的维数，v 是 \(\Phi\)（相应地，Q）的指数，则命题 3 表明有

\[
n \geqslant 2 v.\tag{4}
\]

此外，由于每个全迷向（相应地，全奇异）子空间都包含在某个极大全迷向（相应地，全奇异）子空间之中，故极大全迷向（相应地，全奇异）子空间恰好是维数为 v 的那些。\(\Phi\)（相应地，Q）的指数为 0 这一断言意味着 E 的每个迷向（相应地，奇异）向量都是零。在偶数维 \(n\) 的空间中，中性型恰好是指数为 \(\frac{1}{2} n\) 的那些；在奇数维空间中没有中性型。命题 3 表明，每个型都是一个中性型与一个指数为 0 的型的直和。

命题 4. --- 设 Q 是 E 上的非退化二次型，且存在 E 的向量 \(x \neq 0\) 使得 Q(x) = 0。则对 A 的每个元素 a，存在 \(y \in E\) 使得 Q(y) = a。

事实上，由命题 2 的推论 1，存在子空间 \(\mathbf{G} = \mathbf{F} + \mathbf{F}'\)（\(\mathbf{F}, \mathbf{F}'\)：维数为 1 的全奇异子空间；\(\mathbf{E}\) 的、维数为 2 的），使得 \(\mathbf{Q}\) 在 \(\mathbf{G}\) 上的限制是中性的。若 \(\{e, e'\}\)（\(e \in \mathbf{F}, e' \in \mathbf{F}'\)）是 \(\mathbf{G}\) 的基，则有

\[
\mathrm{Q} (x e + x ^ {\prime} e ^ {\prime}) = b x x ^ {\prime} \quad (x \in \mathrm{A}, x ^ {\prime} \in \mathrm{A}, b \in \mathrm{A}, b \neq 0).
\]

因此取 \(ae + b^{-1}e'\) 为 y 即可。
